% Options for packages loaded elsewhere
\PassOptionsToPackage{unicode}{hyperref}
\PassOptionsToPackage{hyphens}{url}
\PassOptionsToPackage{dvipsnames,svgnames,x11names}{xcolor}
%
\documentclass[
]{scrreprt}

\usepackage{amsmath,amssymb}
\usepackage{iftex}
\ifPDFTeX
  \usepackage[T1]{fontenc}
  \usepackage[utf8]{inputenc}
  \usepackage{textcomp} % provide euro and other symbols
\else % if luatex or xetex
  \usepackage{unicode-math}
  \defaultfontfeatures{Scale=MatchLowercase}
  \defaultfontfeatures[\rmfamily]{Ligatures=TeX,Scale=1}
\fi
\usepackage{lmodern}
\ifPDFTeX\else  
    % xetex/luatex font selection
    \setmainfont[]{Avenir}
    \setsansfont[]{Arial Rounded MT Bold}
\fi
% Use upquote if available, for straight quotes in verbatim environments
\IfFileExists{upquote.sty}{\usepackage{upquote}}{}
\IfFileExists{microtype.sty}{% use microtype if available
  \usepackage[]{microtype}
  \UseMicrotypeSet[protrusion]{basicmath} % disable protrusion for tt fonts
}{}
\makeatletter
\@ifundefined{KOMAClassName}{% if non-KOMA class
  \IfFileExists{parskip.sty}{%
    \usepackage{parskip}
  }{% else
    \setlength{\parindent}{0pt}
    \setlength{\parskip}{6pt plus 2pt minus 1pt}}
}{% if KOMA class
  \KOMAoptions{parskip=half}}
\makeatother
\usepackage{xcolor}
\setlength{\emergencystretch}{3em} % prevent overfull lines
\setcounter{secnumdepth}{-\maxdimen} % remove section numbering
% Make \paragraph and \subparagraph free-standing
\makeatletter
\ifx\paragraph\undefined\else
  \let\oldparagraph\paragraph
  \renewcommand{\paragraph}{
    \@ifstar
      \xxxParagraphStar
      \xxxParagraphNoStar
  }
  \newcommand{\xxxParagraphStar}[1]{\oldparagraph*{#1}\mbox{}}
  \newcommand{\xxxParagraphNoStar}[1]{\oldparagraph{#1}\mbox{}}
\fi
\ifx\subparagraph\undefined\else
  \let\oldsubparagraph\subparagraph
  \renewcommand{\subparagraph}{
    \@ifstar
      \xxxSubParagraphStar
      \xxxSubParagraphNoStar
  }
  \newcommand{\xxxSubParagraphStar}[1]{\oldsubparagraph*{#1}\mbox{}}
  \newcommand{\xxxSubParagraphNoStar}[1]{\oldsubparagraph{#1}\mbox{}}
\fi
\makeatother


\providecommand{\tightlist}{%
  \setlength{\itemsep}{0pt}\setlength{\parskip}{0pt}}\usepackage{longtable,booktabs,array}
\usepackage{calc} % for calculating minipage widths
% Correct order of tables after \paragraph or \subparagraph
\usepackage{etoolbox}
\makeatletter
\patchcmd\longtable{\par}{\if@noskipsec\mbox{}\fi\par}{}{}
\makeatother
% Allow footnotes in longtable head/foot
\IfFileExists{footnotehyper.sty}{\usepackage{footnotehyper}}{\usepackage{footnote}}
\makesavenoteenv{longtable}
\usepackage{graphicx}
\makeatletter
\newsavebox\pandoc@box
\newcommand*\pandocbounded[1]{% scales image to fit in text height/width
  \sbox\pandoc@box{#1}%
  \Gscale@div\@tempa{\textheight}{\dimexpr\ht\pandoc@box+\dp\pandoc@box\relax}%
  \Gscale@div\@tempb{\linewidth}{\wd\pandoc@box}%
  \ifdim\@tempb\p@<\@tempa\p@\let\@tempa\@tempb\fi% select the smaller of both
  \ifdim\@tempa\p@<\p@\scalebox{\@tempa}{\usebox\pandoc@box}%
  \else\usebox{\pandoc@box}%
  \fi%
}
% Set default figure placement to htbp
\def\fps@figure{htbp}
\makeatother

\usepackage{fancyhdr}
\usepackage{hyperref}
\usepackage{wrapfig}
\usepackage{geometry}
\geometry{a4paper, total={170mm,257mm}, left=20mm, top=25mm, bottom=25mm, right=20mm}
\usepackage{lastpage}
\pagestyle{fancy}
\fancyhf{}
\lhead{\textsc{\leftmark}}
\lfoot{}
\rfoot{Page \thepage\ of \pageref*{LastPage}}
\renewcommand{\headrulewidth}{1pt}
\renewcommand{\footrulewidth}{1pt}

% \usepackage[extrabold]{raleway}
% \usepackage{ebgaramond} % Serif font package: EB Garamond
% \usepackage{titlesec}
% \titleformat{\chapter}
% [display]
% {\ralewayextra\Huge\bfseries}
% {\huge\chaptertitlename\space\thechapter}
% {20pt}
% {}

% \titleformat{\section}{\ralewayextra\Large\bfseries}
% \titleformat{\subsection}{\ralewayextra\large\bfseries}

\usepackage{tikz}
\usetikzlibrary{math,positioning,shadows.blur} % TikZ sublibraries
\usepackage{xfp}
\usepackage{ifthen}



% Some colors
\definecolor{Rouge}{HTML}{ffb3ba}
\definecolor{Beige}{HTML}{ffdcb3}
\definecolor{Jaune}{HTML}{ffffb3}
\definecolor{Vert}{HTML}{b3ffb3}
\definecolor{Bleu}{HTML}{b3ccff}
\definecolor{Violet}{HTML}{dfb3ff}
\definecolor{Rose}{HTML}{ffb3ff}

% Parameters
\def\NJ{5}	% Number of days (removes days in command Days is n)
\def\hi{8}	% First hour of the day
\def\hf{19}	% Last hour of the day
\def\H{20.0}% Height of the timetable
\def\D{15.0}% Width of the timetable
\def\h{1.0}	% Height of the header row (days row)
\def\d{1.5} % Width of the first column (times column)
% Days of the timetable. Should contain only the ones displayed
\newcommand\Days{\textbf{Mon}, \textbf{Tue}, \textbf{Wed}, \textbf{Thu}, \textbf{Fri}}
% Constants
\def\dd{\fpeval{(\D-\d) / \NJ}}	% Widht of the columns
\def\Dh{\fpeval{\hf - \hi}}		% Total hours per days (Width of columns)
\def\dh{\fpeval{(\H-\h) / \Dh}}	% Heigth of an hour (Heigth of rows)

% Event command:
% Input: Name1, Name2, Day, Strating time, Duration, Location, Color
% Output: Draw box in TikZ
\NewDocumentCommand{\event}{ O{Prénom} O{Nom} O{\D1} O{\hi} O{1} O{Lieu} O{Beige}}{
% Transform sexagesimal in decimal
\def\Hdebut{\fpeval{floor(#4)}};	% Starting hour (integer)
\def\Mdebut{\fpeval{round((#4 - \Hdebut )*100)}};	% Starting minute (integer)
\def\Debutdeci{\fpeval{\Hdebut + (#4 - \Hdebut ) / 0.6}};	% Starting time (rational)

\def\Hduration{\fpeval{floor(#5)}};	% Duration  hour (integer)
\def\Mduration{\fpeval{round((#5-\Hduration)*100)}};	% Duration minute (integer)
\def\Durationdeci{\fpeval{\Hduration + ((#5-\Hduration) / 0.6)}};	% Duration (rational)

\def\Findeci{\fpeval{\Debutdeci + \Durationdeci}}; % Ending time (rational)
\def\Hfin{\fpeval{floor(\Findeci)}};	% Ending hour (integer)
\def\Mfin{\fpeval{round((\Findeci - \Hfin)*60.0)}};	 % Ending minute (integer)

\def\x{\fpeval{\d + ( #3 - 1 ) * \dd}};	% Columns position (start)
\def\xx{\fpeval{\d + #3 * \dd}};		% Columns position (end)
\def\yy{\fpeval{(\hf - \Debutdeci ) * \dh}};	% Row position (start)
\def\y{\fpeval{\yy - \Durationdeci * \dh}};		% Row position (end)

% Display double zero (i.e. 14:00 instead of 14:0 )
\ifthenelse{\equal{\Mdebut}{0}}
	{
	\ifthenelse{\equal{\Mfin}{0}}
    	{\filldraw [fill= #7, draw=black] (\x, \y) rectangle (\xx, \yy) node[pos=.5, text width=5cm, align=center] {\small\textbf{#1 #2} \linebreak \textit{#6}};}
    	{\filldraw [fill= #7, draw=black] (\x, \y) rectangle (\xx, \yy) node[pos=.5, text width=5cm, align=center] {\small\textbf{#1 #2} \linebreak \textit{#6}};}
    }
    {
    \ifthenelse{\equal{\Mfin}{0}}
    	{\filldraw [fill= #7, draw=black] (\x, \y) rectangle (\xx, \yy) node[pos=.5, text width=5cm, align=center] {\small\textbf{#1 #2} \linebreak \textit{#6}};}
    	{\filldraw [fill= #7, draw=black] (\x, \y) rectangle (\xx, \yy) node[pos=.5, text width=5cm, align=center] {\small\textbf{#1 #2} \linebreak \textit{#6}};}
    }
}

\usepackage{wallpaper} % for the background image on title page
\usepackage{makeidx}
\makeindex
\makeatletter
\@ifpackageloaded{caption}{}{\usepackage{caption}}
\AtBeginDocument{%
\ifdefined\contentsname
  \renewcommand*\contentsname{Table of contents}
\else
  \newcommand\contentsname{Table of contents}
\fi
\ifdefined\listfigurename
  \renewcommand*\listfigurename{List of Figures}
\else
  \newcommand\listfigurename{List of Figures}
\fi
\ifdefined\listtablename
  \renewcommand*\listtablename{List of Tables}
\else
  \newcommand\listtablename{List of Tables}
\fi
\ifdefined\figurename
  \renewcommand*\figurename{Figure}
\else
  \newcommand\figurename{Figure}
\fi
\ifdefined\tablename
  \renewcommand*\tablename{Table}
\else
  \newcommand\tablename{Table}
\fi
}
\@ifpackageloaded{float}{}{\usepackage{float}}
\floatstyle{ruled}
\@ifundefined{c@chapter}{\newfloat{codelisting}{h}{lop}}{\newfloat{codelisting}{h}{lop}[chapter]}
\floatname{codelisting}{Listing}
\newcommand*\listoflistings{\listof{codelisting}{List of Listings}}
\makeatother
\makeatletter
\makeatother
\makeatletter
\@ifpackageloaded{caption}{}{\usepackage{caption}}
\@ifpackageloaded{subcaption}{}{\usepackage{subcaption}}
\makeatother

\usepackage{bookmark}

\IfFileExists{xurl.sty}{\usepackage{xurl}}{} % add URL line breaks if available
\urlstyle{same} % disable monospaced font for URLs
\hypersetup{
  pdftitle={Biometrics in the Bush Capital},
  pdfauthor={International Biometrics Society Australasian Region},
  colorlinks=true,
  linkcolor={blue},
  filecolor={Maroon},
  citecolor={Blue},
  urlcolor={Blue},
  pdfcreator={LaTeX via pandoc}}


\title{Biometrics in the Bush Capital}
\usepackage{etoolbox}
\makeatletter
\providecommand{\subtitle}[1]{% add subtitle to \maketitle
  \apptocmd{\@title}{\par {\large #1 \par}}{}{}
}
\makeatother
\subtitle{24-28th November 2025}
\author{International Biometrics Society Australasian Region}
\date{}

\begin{document}


\begin{titlepage}
\ThisULCornerWallPaper{1.1}{assets/front-cover.png}

\centering
\vspace{4\baselineskip}
{\bfseries\fontsize{50}{50}\selectfont Biometrics in the Bush
Capital \par}
\vspace{3\baselineskip}
{\Large 24-28th November 2025 \par}
\includegraphics[height=20mm]{assets/sponsors/logo.jpg}\\
{\large\textsc{International Biometrics Society Australasian
Region} \par}
\vspace{2\baselineskip}
{\large Sponsored by}\\[5mm]
\centering
Platinum\\[3mm]
\includegraphics[height=10mm]{assets/sponsors/ANU\_Primary\_Horizontal\_GoldBlack.png}
\par
Gold\\[3mm]
\includegraphics[height=10mm]{assets/sponsors/GRDC\_Logo\_Primary\_Default\_transparency.png}
\includegraphics[height=10mm]{assets/sponsors/UoA-Logo-Primary-RGB-Large.png}
\par
Supporter\\[3mm]
\includegraphics[height=10mm]{assets/sponsors/SSA\_Logotype\_MAIN
USE.png}
\par

\end{titlepage}

\newpage 

\renewcommand*\contentsname{Table of contents}
{
\hypersetup{linkcolor=}
\setcounter{tocdepth}{1}
\tableofcontents
}

\chapter{About}\label{about}

\begin{wrapfigure}{L}{2.5cm}
\includegraphics[width=2.5cm]{assets/sponsors/logo.jpg}
\end{wrapfigure}

The \href{https://www.biometricsociety.org/}{International Biometric
Society (IBS)} is devoted to the development and application of
statistical and mathematical theory and methods in the Biosciences,
including agriculture, biomedical science and public health, ecology,
environmental sciences, forestry, and allied disciplines. It welcomes as
members statisticians, mathematicians, biological scientists, and others
devoted to interdisciplinary efforts in advancing the collection and
interpretation of information in the biosciences.

The IBS is organised into geographically-defined regions, each with at
least 50 members. \href{https://biometricsociety.org.au/}{The
Australasian Region (IBS-AR)} is one of over 30 regions under the IBS
umbrella. It provides to members of the IBS who reside in Australia, New
Zealand and nearby countries a local focus and support for biometrics
undertaken in this region.

The Australasian Region Biometrics conference is a biannual conference
sponsored by the IBS-AR. The conference is a forum for biometricians,
statisticians, and those interested in the development and application
of statistical and mathematical theory and methods to problems in the
biosciences. The conference is open to members and non-members of the
Society.

\Large Conference website:
\href{https://biometricsociety.org.au/conference2025}{biometricsociety.org.au/conference2025}

\normalsize

\section{IBS-AR Council}\label{ibs-ar-council}

\begin{itemize}
\tightlist
\item
  \textbf{Alison Kelly} (University of Queensland) as President
\item
  \textbf{James Curran} (University of Auckland) as Vice-President
\item
  \textbf{Louise McMillan} (Victoria University of Wellington) as
  Secretary
\item
  \textbf{Ian Renner} (University of Newcastle) as Treasurer
\item
  \textbf{Sam Rogers} (University of Adelaide) as Membership Officer
\item
  \textbf{Emi Tanaka} (Australian National University) as Social Media
  Officer
\item
  \textbf{Warren Muller} (CSIRO) as Public Officer
\item
  \textbf{Charley Budgeon} (University of Western Australia) as
  Biometric Bulletin Correspondent
\item
  \textbf{Peter Green} (AgResearch Ltd) as Website manager
\item
  \textbf{Samuel Muller} (Macquarie University) as IBS Council Member
\item
  \textbf{Kevin Murray} (University of Western Australia) as IBS Council
  Member
\item
  \textbf{Chris Triggs} (University of Auckland) as IBS Executive Board
\end{itemize}

\vfill

\emph{Conference booklet compiled by Emi Tanaka}\\
\emph{Last updated: 2025-09-09 08:08:11}

\chapter{Organisers}\label{organisers}

\section{Local Organising Committee}\label{local-organising-committee}

\begin{itemize}
\tightlist
\item
  \textbf{Emi Tanaka} (Australian National University) as Chair
\item
  \textbf{Alan Welsh} (Australian National University)
\item
  \textbf{Fonti Kar} (Australian National University)
\item
  \textbf{Joanne Thandrayen}
\item
  \textbf{Warren Muller} (CSIRO)
\item
  \textbf{Weihao (Patrick) Li} (Australian National University)
\item
  \textbf{Yidi Deng} (Australian National University)
\end{itemize}

\section{Social Program Committee}\label{social-program-committee}

\begin{itemize}
\tightlist
\item
  \textbf{Quan Vu} (Australian National University) as Chair
\end{itemize}

\section{Scientific and Program
Committee}\label{scientific-and-program-committee}

\begin{itemize}
\tightlist
\item
  \textbf{David Warton} (UNSW Sydney) as Chair
\item
  \textbf{Carolyn Huston} (CSIRO)
\item
  \textbf{Emi Tanaka} (Australian National University)
\item
  \textbf{Jean Yang} (University of Sydney)
\item
  \textbf{Julian Taylor} (Adelaide University)
\item
  \textbf{Nicola Armstrong} (Curtin University)
\item
  \textbf{Thomas Lumley} (University of Auckland)
\end{itemize}

\section{Conference Allies}\label{sec-allies}

\begin{itemize}
\tightlist
\item
  \textbf{Fonti Kar} (Australian National University)
\item
  \textbf{Louise McMillan} (Victoria University of Wellington)
\item
  \textbf{Scott Foster} (CSIRO)
\end{itemize}

\chapter{Sponsors}\label{sponsors}

We gratefully acknowledge our sponsors for their generous support, which
has made this conference possible.

\vfill

\begin{center}

{\Large Platinum  }

\includegraphics[width=50mm]{assets/sponsors/ANU_Primary_Horizontal_GoldBlack.png}  


\vfill
  

{\Large Gold}


\includegraphics[width=50mm]{assets/sponsors/GRDC_Logo_Primary_Default_transparency.png}  \includegraphics[width=50mm]{assets/sponsors/UoA-Logo-Primary-RGB-Large.png}  


\vfill


{\Large Supporters}

\includegraphics[width=50mm]{assets/sponsors/SSA_Logotype_MAIN USE.png}

ACT Branch

\end{center}

\chapter{Schedule}\label{schedule}

\pagestyle{empty}

\begin{tikzpicture}

\draw [draw=black] (0.0,0.0) rectangle (15.0,20.0);

\def\fin{\fpeval{\Dh-1}};
\foreach \i in {1,...,\fin}{
\def\y{\fpeval{\i*\dh}} ;
\def\heure{\fpeval{(\hf-\i)}};
\def\x{\fpeval{\d/2.0}}
\node at (\x,\y)  {\heure: 00};
\draw[dashed](\d, \y)--(\D, \y);
};

\foreach \j [count=\i] in \Days{
\def\y{\fpeval{\H-0.5*\h}};
\def\x{\fpeval{2*\d + (\i - 1) * \dd - 0.05}};
\node at (\x , \y)  {\j};
};

\event[Registration][][1][8.30][0.30][Foyer];
\event[Registration][][2][8.30][0.30][Foyer];
\event[Registration][][3][8.30][0.30][Foyer];
\event[Registration][][4][8.30][0.30][Foyer];
\event[Registration][][5][8.30][0.30][Foyer];

\event[Workshop][][1][9.00][8][];


\end{tikzpicture}

\pagestyle{fancy}

\chapter{Keynote talks}\label{keynote-talks}

\subsection{Optimizing research impact through interdisciplinary and
collaborative
research}\label{optimizing-research-impact-through-interdisciplinary-and-collaborative-research}

\emph{Charmaine B. Dean}\index{Dean, Charmaine B.} \emph{(University of
Waterloo)}

\setlength{\parskip}{0.5em}

Interdisciplinary collaborative research is a key component of data
science and for some of us, plays an important part of our roles as
statisticians. It is not unusual that we become accustomed to vertical
thinking whereby we use existing tools and methods in our own specialty
to problem solve, losing sight of the larger interdisciplinary context
of data science, and the context of the scientific challenge. The
Government of Canada - Science and Technology branch has identified
several priority research challenge topics that involve
cross-disciplinary work. Although statistical tools and analytics are
identified in these research challenge priority areas, additionally, the
development of fundamental transformative and enabling technological
tools specifically for statistical methods and analytics to support
research and societal advancement is also seen as a priority. This talk
shares insights about the challenges and opportunities for statistics in
interdisciplinary research. Specifically, monitoring viral signals in
wastewater and assessing forest fire risk are given as complex, case
studies that use a collaborative and interdisciplinary approach to solve
difficult problems. This approach will demonstrate the significant
benefits for not only optimizing research impact but for training
students to become horizontal problem solvers across a wide range of
research methods which will benefit them in navigating complex problems
and in the development of appropriate tools for their analysis

\subsection{Saddlepoint approximations for
likelihoods}\label{saddlepoint-approximations-for-likelihoods}

\emph{Jesse Goodman}\index{Goodman, Jesse} \emph{(University of
Auckland)}

\emph{Co-authors: Godrick Oketch and Rachel Fewster; Alice Hankin,
Sixiang Shan and George Fan}

\setlength{\parskip}{0.5em}

Classically, the saddlepoint approximation has been used as a systematic
method for converting a known generating function into an approximation
for an unknown density function. More recently, it has been used instead
as an approximation to the \emph{likelihood function}. In this
viewpoint, it is the underlying data-generating process whose generating
function is used, and the saddlepoint approximation can be maximized to
compute an approximate saddlepoint MLE for given observed data. This
talk will explain how the saddlepoint approximation can be interpreted
with a statistical lens, including common features for those otherwise
intractable models where we can compute a generating function but not a
likelihood. Many of these models come from statistical ecology,
including sitations where we gather population-wide observations only.
In addition, the talk will describe a class of models having simple
theoretical guarantees for the effect of using the saddlepoint MLE as a
substitute for the unknown true MLE, and will demonstrate new tools to
visualize the saddlepoint approximation intuitively, to simplify and
automate the computation of saddlepoint MLEs, and to quantitatively
assess the amount of approximation error introduced by using an
approximate likelihood as a substitute for an intractable true
likelihood.

\subsection{Modularizing Biometric Models Facilitates Multistage
Computing}\label{modularizing-biometric-models-facilitates-multistage-computing}

\emph{Mevin B. Hooten}\index{Hooten, Mevin B.} \emph{(The University of
Texas at Austin)}

\setlength{\parskip}{0.5em}

Bayesian modeling has become invaluable in biometrics. It allows us to
formally consider unobserved processes while accommodating uncertainty
about data collection and our understanding of biological and ecological
mechanisms. Several excellent software packages are available for
fitting Bayesian models to data and are being applied every day to
analyze biometric data. These methods allow us to answer questions using
data in ways that has never before been possible. The adoption of
Bayesian methods has led to bigger models necessary to answer tough
questions using large and varied data sets. Bigger models and data sets
lead to computing bottlenecks. Fortunately, a solution to Bayesian
computing roadblocks sits in plain sight. The structure of Bayesian
models allows us to rearrange them so that we can perform computing in
stages. We can break big models into pieces, fit them separately, and
then recombine them in later computing stages. Recursive Bayesian
approaches can save us time by leveraging the parallel architecture of
modern computers. A modular perspective allows us to see Bayesian models
in a way that facilitates multistage computing. I will demonstrate the
procedure with a set of biometric examples. These include geostatistical
models in marine science, capture-recapture models for abundance
estimation, and spatial point process models for species distributions.

\subsection{On Finding Good
Experiments}\label{on-finding-good-experiments}

\emph{Cheng Soon Ong}\index{Ong, Cheng Soon} \emph{(CSIRO and ANU)}

\setlength{\parskip}{0.5em}

One of the key choices we have as scientists is to design informative
experiments. With computational methods like AI promising accurate
predictions, we revisit the question of adaptively designing new
measurements that take previous data into account. Using examples from
genomics, we illustrate some recent ideas on using machine learning to
recommend experiments. Then we discuss potential impacts on choosing
measurements in spatiotemporal problems. We conclude by outlining some
opportunities and challenges of including machine learning in the
scientific discovery process.

\subsection{Uses of gnm for Generalized (Non)linear
Modelling}\label{uses-of-gnm-for-generalized-nonlinear-modelling}

\emph{Heather L. Turner}\index{Turner, Heather L.} \emph{(University of
Warwick)}

\setlength{\parskip}{0.5em}

The R package \{gnm\} was designed as a unified interface to fit
Generalized Nonlinear Models: \emph{generalized} to handle responses
with restricted range and/or a variance that depends on the mean, and
\emph{nonlinear} to allow the predictor for the mean to be nonlinear in
its parameters. This framework covers several models that were proposed
in the literature and adopted in practice before \{gnm\} was released,
but used to require a mixed bag of specialised software to fit.

With \{gnm\} celebrating its 20th birthday this year, it's a good time
to review how the package is being used. I'll highlight some of the
applications we were aware of when \{gnm\} was first developed, that
remain in common use, and explore more recent applications, particularly
in the field of biometrics.

We'll discover one motivation for using \{gnm\}, is for the
``eliminate'' feature that efficiently estimates stratification
parameters. This can be useful even when the predictor is linear, as in
the case of using conditional Poisson models to analyse case-crossover
studies in epidemiology. We'll also look at two of the packages that
have built on \{gnm\}. The first, \{multgee\}, uses \{gnm\} to fit
multiplicative interactions for certain correlation structures when
modelling categorical data, with applications in public health,
agriculture, and psychology. The second, \{VFP\}, is a more specialised
package that uses \{gnm\} to model the mean-variance relationship for
in-vitro diagnostic assays.

Through these use cases we'll see how different features of \{gnm\} can
be applied, demonstrating the versatility of this software.

\chapter{Invited talks}\label{invited-talks}

\subsection{Inferring the rate of undetected contamination using random
effects modelling of biosecurity screening
histories}\label{inferring-the-rate-of-undetected-contamination-using-random-effects-modelling-of-biosecurity-screening-histories}

\emph{Sumonkanti Das}\index{Das, Sumonkanti} \emph{(Australian National
University)}

\setlength{\parskip}{0.5em}

Group testing plays a vital role in biosecurity operations worldwide,
particularly in minimising the risk of introducing exotic pests,
contaminants, and pathogens through imported agricultural products. A
common screening strategy involves pooling items from consignments and
testing each group for contamination presence, with consignments
typically rejected if any group tests positive. Although screening
designs often target a high probability of detection assuming a fixed
minimum prevalence, analysing the historical results of these tests to
infer the extent of contamination in non-rejected consignments (referred
to as leakage) is less common.

This study advances censored beta-binomial (BB) models to address
contamination risk in frozen seafood imports into Australia,
incorporating imperfect tests. Motivated by the characteristics of our
case study, we develop a new class of BB models that impose a minimum
positive consignment propensity threshold, capturing scenarios where
contamination is either absent or exceeds a known minimum level. To fit
these models, we propose a Metropolis-Hastings (MH) algorithm
conditioned on prior distributions for sensitivity and specificity,
allowing efficient estimation of quantities related to contamination
levels. We analyse historical testing data under multiple scenarios
using the proposed MH algorithm, yielding novel insights into both
contamination risk and leakage.

Finally, we use model-based simulations to communicate risk levels,
providing key insights into potential undetected contamination.

\subsection{Predication of Daily Weight Gain with Cattle Behaviour and
Daily Activity Using Triaxial Accelerometer
Data}\label{predication-of-daily-weight-gain-with-cattle-behaviour-and-daily-activity-using-triaxial-accelerometer-data}

\emph{Shuwen Hu}\index{Hu, Shuwen} \emph{(RMIT)}

\emph{Co-authors: Reza Arablouei, Greg Bishop-Hurley, Neil Bagnall and
Aaron Ingham}

\setlength{\parskip}{0.5em}

Live body weight gain is an important measurement of animal performance.
In this study, we predict the cattle's daily weight gain from five core
cattle behaviours and a measure of total daily activity based on
accelerometer data. To collect data, we conducted an experiment
equipping a herd of 60 Brahman steers with research collars containing
triaxial accelerometers over nearly one month in Australia. We used the
accelerometer data, which represents the intensity of animal movement,
to compute an activity metric within a five-minute window. In addition,
we use pre-trained accelerometer-based machine learning models to
classify cattle behaviour into grazing, ruminating, resting, walking,
drinking or other classes over five-second time windows. Daily behaviour
profiles were constructed for each animal and experiment day by
aggregating the behaviour predictions over every calendar day. Our
objective was to explore how to use behaviours and activity metrics to
predict the cattle's daily weight gain. The daily activity values
ranging from 5.44g to 23.69g. The average daily time spent grazing,
ruminating, resting, walking and drinking was 8.97±1.12, 7.78 ±1.03,
5.83±1.05, 1.00±0.4, and 0.18±0.12 hours, respectively. Some weather
information data are combined in the model to predict the cattle
live-weight gain. The best R-squared value is 0.467, with a minimum root
mean square error (RMSE) of 0.867 from the linear regression model. The
live-weight gain could not be fully explained by measurements taken in
this study, but we showed how these factors can influence the
variability in cattle performance.

\subsection{The difficulties of clustering categorical or mixed
data}\label{the-difficulties-of-clustering-categorical-or-mixed-data}

\emph{Louise McMillan}\index{McMillan, Louise} \emph{(Victoria
University of Wellington)}

\emph{Co-authors: Daniel Fernández, Shirley Pledger, Richard Arnold, Ivy
Liu, Eleni Matechou, Adam Glucksman}

\setlength{\parskip}{0.5em}

I will discuss techniques for clustering categorical data that go beyond
treating the data as numerical or converting the categories to dummy
variables. I will talk about a range of approaches to clustering
categorical data, including the \texttt{clustord} R package for
clustering ordinal data, which uses model-based clustering via
likelihoods and the EM algorithm. Then I will discuss a Bayesian
approach from population genetics that may be extendable to general
mixed datasets and is the subject of my latest research.

\subsection{Evaluating the impact of trait measurement error on genetic
analysis of computer vision-based
phenotypes}\label{evaluating-the-impact-of-trait-measurement-error-on-genetic-analysis-of-computer-vision-based-phenotypes}

\emph{Gota Morota}\index{Morota, Gota} \emph{(The University of Tokyo)}

\setlength{\parskip}{0.5em}

Quantitative genetic analysis of image- or video-derived phenotypes is
increasingly being performed for a wide range of traits. Pig body weight
values estimated by a conventional approach or a computer vision system
can be considered as two different measurements of the same trait, but
with different sources of phenotyping error. Previous studies have shown
that trait measurement error, defined as the difference between manually
collected phenotypes and image-derived phenotypes, can be influenced by
genetics, suggesting that the error is systematic rather than random and
is more likely to lead to misleading quantitative genetic analysis
results. Therefore, we investigated the effect of trait measurement
error on genetic analysis of pig body weight (BW). Calibrated
scale-based and image-based BW showed high coefficients of determination
and goodness of fit. Genomic heritability estimates for scale-based and
image-based BW were mostly identical across growth periods. Genomic
heritability estimates for trait measurement error were consistently
negligible, regardless of the choice of computer vision algorithm. In
addition, genome-wide association analysis revealed no overlap between
the top markers identified for scale-based BW and those associated with
trait measurement error. Overall, the deep learning-based regressions
outperformed the adaptive thresholding segmentation methods. This study
showed that manually measured scale-based and image-based BW phenotypes
yielded the same quantitative genetic results. We found no evidence that
BW trait measurement error could be influenced, at least in part, by
genetic factors. This suggests that trait measurement error in pig BW
does not contain systematic errors that could bias downstream genetic
analysis.

\subsection{Optimal allocation of resources between control and
surveillance for complex eradication
scenarios}\label{optimal-allocation-of-resources-between-control-and-surveillance-for-complex-eradication-scenarios}

\emph{Mahdi Parsa}\index{Parsa, Mahdi} \emph{(Department of Agriculture,
Fisheries and Forestry)}

\setlength{\parskip}{0.5em}

Effective eradication of invasive species over large areas requires
strategic allocation of resources between control measures and
surveillance activities. This study presents an analytical Bayesian
framework that integrates stochastic modelling and explicit measures of
uncertainty to guide decisions in complex eradication scenarios. By
applying Shannon entropy to quantify uncertainty and incorporating the
expected value of perfect information (EVPI), the framework identifies
conditions under which investment into control or surveillance becomes
worthwhile. Findings show that strategies which hedge against
uncertainty can markedly improve the robustness of eradication outcomes
with only marginal increases in expected costs. This approach offers
practical tools for designing more cost-effective and reliable
eradication programs and for prioritising data collection to reduce
uncertainty where it has the greatest impact.

\subsection{Surveillance for counterfactual scenarios of invasive
species - why it's useful and a convenient ~way to do
it.}\label{surveillance-for-counterfactual-scenarios-of-invasive-species---why-its-useful-and-a-convenient-way-to-do-it.}

\emph{Andrew P. Robinson}\index{Robinson, Andrew P.} \emph{(The
University of Melbourne)}

\emph{Co-authors: John Kean, Melissa Welsh, and Kuo-Szu Chiang}

\setlength{\parskip}{0.5em}

The impact of invasive species is affected by a range of factors, many
of which can be anticipated in advance -- for example, the prevalence of
host material, climate suitability, the size of affected agricultural
resources, and so on. One factor that cannot be anticipated in advance
is the size of the incursion at the time of its detection.
Unfortunately, the impact of an incursion is tightly tied to its
maturity at detection, ranging from a single seed, for example, to a
50,000 hectare infestation of plants.

We propose a simple probability model for the detection of an invasive
species that can either capture or integrate out the consequent
uncertainty of the maturity of the incursion. We represent the
relationship between surveillance and the detection of the organism
using survival analysis: the detection of the incursion is analogous to
the survival event; it is a binary occurrence that happens at some point
in time, and once it has happened it does not happen again.

Under such a model, we can connect the distribution of the size of the
infestation at the time to detection to the probability of detecting the
incursion given that it has not already been detected, namely, the
hazard function. For example, a popular model for the detection of an
incursion of size x with number of traps t and probability of detecting
a single pest \(p\) is:

\(h(x,t,p) = 1 - (1 - p)^{tx}\)

Algebra leads us to a size-at-detection pdf. Other corrections are also
applied as needed. The outcome is a pdf that is a function of process
parameters, enabling straightforward assessment of different
surveillance choices. Parameter estimates for the distribution can be
derived from first principles, field experiments, or expert elicitation.

In this presentation we will derive and demonstrate the use of the
survival-based incursion size at detection pdf and discuss its
implications and challenges.

\subsection{Fishing for Heritability in the Gill Microbiome: Why
Statisticians Should get out into the
field}\label{fishing-for-heritability-in-the-gill-microbiome-why-statisticians-should-get-out-into-the-field}

\emph{Elle Saber}\index{Saber, Elle} \emph{(Australian National
University)}

\setlength{\parskip}{0.5em}

Host‐associated microbiomes are increasingly recognised as integral to
health, yet the extent to which host genetics shapes these communities
remains unclear. While heritable components of gut and skin microbiomes
have been documented in several vertebrates, evidence for the gill
microbiome of fish is scarce. This exploratory study sampled the gill
microbiome of Atlantic salmon within a Tasmanian breeding program to
investigate potential genetic influences. Despite careful planning,
study participants did not always behave, and the practical constraints
of a commercial operation prevented complete data capture. Having
participated in the data collection I was better prepared to understand
the limitations of the dataset when doing the downstream analysis. The
take home message is that, even if we feel like a fish out of water,
time spent among the data can be just as valuable as time spent our
desks.

\subsection{Bayesian clustered ensemble prediction for multivariate time
series}\label{bayesian-clustered-ensemble-prediction-for-multivariate-time-series}

\emph{Shonosuke Sugasawa}\index{Sugasawa, Shonosuke} \emph{(Keio
University)}

\setlength{\parskip}{0.5em}

We propose a novel methodology called the mixture of Bayesian predictive
syntheses (MBPS) for multiple time series count data and apply the
methodology to predict the numbers of COVID-19 inpatients and isolated
cases in Japan and Korea at the subnational level. MBPS combines a set
of predictive models and partitions the multiple time series into
clusters based on their contribution to predicting the outcome. In this
way MBPS leverages the shared information within each cluster and avoids
using a multivariate count model, which is generally cumbersome to
develop and implement. Our data analyses demonstrate that the proposed
MBPS methodology has improved predictive accuracy and uncertainty
quantification.

\subsection{Optimal sampling in border biosecurity: Application to
skip-lot
sampling}\label{optimal-sampling-in-border-biosecurity-application-to-skip-lot-sampling}

\emph{Raphael Trouve}\index{Trouve, Raphael} \emph{(The University of
Melbourne)}

\setlength{\parskip}{0.5em}

Border biosecurity faces mounting pressure from increasing global trade,
requiring cost-effective inspection strategies to reduce the risk of
importing pest and diseases. Current international standards recommend
inspecting all incoming consignments (full census) with fixed sample
sizes (e.g., 600 units) for high-risk pathways, but this may be overkill
for lower-risk pathways with established compliance records. When should
agencies use skip-lot sampling (SLS, sometimes called continuous
sampling plan), which adaptively reduces inspections based on recent
compliance history, over full census inspection?

We developed a propagule pressure equation for SLS in overdispersed
pathways and used Lagrange multipliers to derive a solution. Results
show the choice depends on pathway overdispersion, sampling costs, and
budget constraints. Optimal sample sizes are typically smaller than
current recommendations, with better returns from inspecting a larger
proportion of consignments rather than larger samples per consignment.
This framework provides biosecurity agencies with data-driven guidance
for implementing adaptive sampling strategies.

\subsection{Species archetype models for presence-only
data}\label{species-archetype-models-for-presence-only-data}

\emph{Skipton N.C. Woolley}\index{Woolley, Skipton N.C.} \emph{(CSIRO)}

\emph{Co-authors: David Warton, Piers Dunstan}

\setlength{\parskip}{0.5em}

Joint species distribution modelling is a recently emerged and
potentially powerful statistical method for analysing biodiversity data.
Despite the plethora of presence-only occurrence data available for
biodiversity analysis, there remain few examples of presence-only
multiple-species modelling approaches. We present a
mixture-of-regressions model for understanding how groups of species are
distributed based on presence-only data. Our approach extends Species
Archetype Models using a point process framework and incorporates joint
estimation of sighting bias based on all species occurrences included in
the model. We demonstrate that our method can accurately recapture mean
and variance of parameters from simulated data sets and provide better
estimates than those generated from multiple single species
presence-only species distribution models. We apply our approach to a
Myrtaceae presence-only occurrence dataset from New South Wales,
Australia. We show that presence-only Species Archetype Models allow for
the propagation of variance and uncertainty from the data through
predictions, improving inference made on presence-only biodiversity data
for multiple species.

\chapter{Contributed talks}\label{contributed-talks}

\subsection{Handling Missingness in Prevalence Estimates from National
Surveys}\label{handling-missingness-in-prevalence-estimates-from-national-surveys}

\emph{Oyelola Adegboye}\index{Adegboye, Oyelola} \emph{(Menzies School
of Health Research)}

\emph{Co-authors: Li Siyu, Tomoki Fujii and Denis HY Leung}

\setlength{\parskip}{0.5em}

Population-based surveys, such as the Demographic and Health Surveys
(DHS), are pivotal for estimating the prevalence of important diseases,
particularly in low-resource settings. However, missing data, such as
person non-response, particularly refusals, poses a serious challenge,
potentially introducing substantial bias in prevalence estimates. Using
HIV data sets from Malawi's 2004 DHS, Antenatal Clinics surveillance,
and the Malawi Diffusion and Ideational Change Project (MDICP), this
paper evaluates existing estimators and proposes novel approaches to
adjust for refusal bias. These include complete case analysis, mean
score imputation, inverse propensity score weighting, bounding
techniques using longitudinal or sentinel data and Manski’s partial
identification approach. The study distinguishes between refusals and
non-contacts and examines how prior knowledge of HIV status influences
participation. Estimates of HIV prevalence varied notably across
methods, with refusal-adjusted approaches generally yielding higher
prevalence rates than complete case estimates. The divergence is more
pronounced among women, largely due to higher refusal rates among them.
Bounding methods provide credible intervals for prevalence under weak
assumptions. In the case of HIV estimates, refusal bias, particularly
when linked to prior testing knowledge, can significantly distort HIV
estimates. Integrating multiple data sources and using methodologically
transparent adjustment techniques are critical for robust HIV
surveillance in low-resource settings.

\subsection{Reporting Odds Ratios under Fluctuating Reporting Rates in
Spontaneous Reporting
Systems}\label{reporting-odds-ratios-under-fluctuating-reporting-rates-in-spontaneous-reporting-systems}

\emph{Tatsuhiko Anzai}\index{Anzai, Tatsuhiko} \emph{(Institute of
Science Tokyo)}

\emph{Co-authors: Kunihiko Takahashi}

\setlength{\parskip}{0.5em}

Spontaneous adverse event reporting systems, including the Japanese
Adverse Drug Event Report (JADER) database, the FDA Adverse Event
Reporting System (FAERS), and VigiBase, play a critical role in
post-marketing drug safety surveillance. The reporting odds ratio (ROR)
is a commonly used measure for detecting suspected adverse drug
reactions as ``signals'', based on disproportionality in reporting
between a specific drug and all others, as summarized in a two-by-two
contingency table. However, during events such as the COVID-19 pandemic,
reporting rates of adverse drug reactions can fluctuate, potentially
introducing bias into ROR. This study proposes a method for deriving the
ROR that incorporates four parameters, each corresponding to
fluctuations in reporting ratios for the cells in the two-by-two
contingency table. These parameters can be estimated using the
divergence between observed and predicted values, where the predicted
values are derived from a regression model under the assumption that the
reporting rates for the drug and the adverse reaction remain stable. We
evaluate the properties and performance of the proposed method through
both real-world data analysis and simulation studies, demonstrating its
effectiveness for signal detection.

\subsection{Pooled testing with penalized regression
models}\label{pooled-testing-with-penalized-regression-models}

\emph{Christopher Bilder}\index{Bilder, Christopher} \emph{(University
of Nebraska-Lincoln, Department of Statistics)}

\emph{Co-authors: Pranta Das, Joshua Tebbs, and Christopher McMahan}

\setlength{\parskip}{0.5em}

Pooled testing (also known as group testing) is a widely used procedure
to test individuals for infectious diseases. Rather than testing each
specimen separately, multiple specimens are pooled together and tested
as one. The pool test outcome, along with further tests as needed, are
used to determine which individuals are positive or negative for an
infectious pathogen. The COVID-19 pandemic especially highlighted the
importance of pooled testing, with laboratories adopting it worldwide to
increase their testing capacity. Pooled testing traditionally relies on
observing the positive/negative test outcomes alone. However, during the
pandemic, new pooled testing algorithms were developed that utilize the
viral load information from a pool test. These new algorithms use
penalized regression models to predict the viral load of each
individual, which lead to individual positive/negative predictions. The
purpose of our presentation is to provide a comparison of these new
algorithms relative to standard ones. Both gains and losses are
quantified by applying algorithms under fair comparison settings.

\subsection{Nested-factorial treatment models: types, their uses and
examples}\label{nested-factorial-treatment-models-types-their-uses-and-examples}

\emph{Chris Brien}\index{Brien, Chris} \emph{(The Plant Accelerator,
University of Adelaide)}

\setlength{\parskip}{0.5em}

Most commonly factorial experiments involve several crossed treatment
factors and the treatment model includes terms for all possible
combinations of the treatment factors. However, treatment models that
employ nested treatment factors are beneficial in a number of situations
and could be used more often. Applicable situations include those in
which (i) treatment factors are intrinsically nested, (ii) treatments
involve one or more factors plus a control and (iii) treatments have
been systematically allocated and/or are pseudoreplicated. Four basic
types of nested factors that start with a single nesting factor are (i)
one nested factor, (ii) multiple, crossed, nested factors, (iii)
multiple, independent, nested factors, and (iv) a hierarchy of multiple,
nested factors. Examples involving the different situations and types of
nesting will be presented. Constructing the factors using the \texttt{R}
package \texttt{dae} (Brien, 2025,
\url{https://cran.r-project.org/package=dae}) will be described and the
properties of the models discussed for the examples via their anatomies
produced using \texttt{dae}.

\subsection{Bayesian Ordinal Regression for Crop Development and Disease
Assessment}\label{bayesian-ordinal-regression-for-crop-development-and-disease-assessment}

\emph{Zhanglong Cao}\index{Cao, Zhanglong} \emph{(Centre for Crop and
Disease Management, Curtin University)}

\emph{Co-authors: Rose Megirian, Matthew Nguyen and Adam Sparks}

\setlength{\parskip}{0.5em}

Accurate assessment of crop development and disease severity is
essential for informed agronomic decision-making. This study presents a
Bayesian framework for analysing ordinal data from field trials,
including growth scale progression and disease ranking scores. Using the
brms package in R, we applied cumulative logit models to evaluate the
effects of sowing depth and treatment combinations on cereal growth
stages, measured via the Zadok’s scale, across two Western Australian
sites (Merredin and Wickepin). The same framework is being extended to
model disease severity scores, demonstrating its versatility across
categorical biological measurements. Our workflow incorporates rigorous
model testing and evaluation, including posterior predictive checks and
leave-one-out cross-validation (LOO-CV), to ensure robust inference and
model fit. Rather than relying on p-values from linear mixed models, the
Bayesian approach provides interpretable probabilities of achieving
specific growth stages or disease severity levels. This shift enables
more nuanced understanding of treatment effects and supports
decision-making under uncertainty.

\subsection{Do Mice Matter? The Impact of Mice on a New Zealand
Ecosanctuary}\label{do-mice-matter-the-impact-of-mice-on-a-new-zealand-ecosanctuary}

\emph{Vanessa Cave}\index{Cave, Vanessa} \emph{(VSNi)}

\setlength{\parskip}{0.5em}

Mammal-resistant fences have enabled the successful eradication of
exotic mammals from ecosanctuaries in New Zealand. However, preventing
the re-invasion of mice remains a challenge. Indeed, mice are still
present in many fenced ecosanctuaries and can reach high population
densities.

Scientists at Manaaki Whenua – Landcare Research have been studying
the impact of mice on biodiversity at Sanctuary Mountain Maungatautari.
Two independently fenced sites within the sanctuary were managed to
create contrasting mouse populations: one with high mouse numbers, and
the other with undetectable levels. After two years, management
protocols were reversed, with mice eradicated from the first site and
allowed to increase at the second. Data on the abundance of
invertebrates, seedlings, and fungi were collected throughout the
duration of the study.

Temporal trends in abundance were analysed and compared using linear
mixed models with smoothing splines. The findings suggest that mice can
have severe impacts on native biodiversity, particularly invertebrates,
posing a significant threat to ecological recovery efforts in fenced
ecosanctuaries.

\subsection{Automatic debiased machine learning (autoDML) for causal
inference: implementation and evaluation in real-world observational
studies}\label{automatic-debiased-machine-learning-autodml-for-causal-inference-implementation-and-evaluation-in-real-world-observational-studies}

\emph{Tong Chen}\index{Chen, Tong} \emph{(Murdoch Children's Research
Institute)}

\emph{Co-authors: Stijn Vansteelandt and Margarita Moreno-Betancur}

\setlength{\parskip}{0.5em}

The estimation of average treatment effect (ATE) is a central goal of
observational health research. Causal machine learning (CML) methods are
increasingly popular for ATE estimation, as they leverage data-adaptive
algorithms to account for complex relationships amongst confounders,
exposure and outcome while ensuring valid inference. Major CML methods,
such as double machine learning and targeted maximum likelihood
estimation, offer robustness and efficiency advantages but they require
first estimating the propensity score and then calculating treatment
divided by the propensity score, known as the Riesz representer.
Standard machine learning algorithms for the propensity score estimation
are suboptimal as they are designed to minimise prediction error,
ignoring that equally large estimation errors in small and large
propensity scores propagate very differently into the ATE estimates upon
taking the reciprocal. Automatic debiased machine learning (autoDML),
introduced by Chernozhukov et al.~(2021), addresses this issue by
directly estimating the Riesz representer using machine learning
methods, yielding a more stable alternative to standard CML methods.
Despite its theoretical advances, autoDML remains largely unknown in
biostatistical research and has not been compared to standard CML
methods in realistic simulation studies. In this talk, we bridge these
gaps by introducing autoDML and highlighting its advantages. We describe
a practical, step-by-step guide for implementing autoDML, and compare
autoDML with standard CML methods using comprehensive, realistic
simulation studies. We further illustrate the utility of autoDML using
data from the Longitudinal Study of Australian Children to evaluate the
impact of overweight or obesity on cardiovascular outcomes in
adolescence.

\subsection{A Set of Precise Asymptotics for Gaussian Variational
Approximation in Generalised Linear Mixed
Models}\label{a-set-of-precise-asymptotics-for-gaussian-variational-approximation-in-generalised-linear-mixed-models}

\emph{Nelson J. Y. Chua}\index{Chua, Nelson J. Y.} \emph{(Australian
National University)}

\emph{Co-authors: Francis K. C. Hui, Luca Maestrini, A. H. Welsh}

\setlength{\parskip}{0.5em}

Generalised linear mixed models (GLMMs) are used to model data with a
clustered or hierarchical structure, capturing intra-cluster dependence
through the use of (latent) random effects. However, this dependence
structure results in a likelihood function that is typically
computationally intensive to evaluate, and so approximate likelihood
approaches are often used instead for model fitting and inference. One
such approach which has grown in popularity over the past decade is
Gaussian variational approximation (GVA). In addition to estimates of
the model parameters, GVA concurrently produces random effects
predictions and associated uncertainty measures.

In this talk, we formulate a set of precise asymptotic properties for
both the parameter estimates and random effects predictions obtained
from GVA, focusing on independent cluster GLMMs. We find that these
properties change substantially depending on whether or not the random
effects are conditioned on. By comparing GVA's asymptotic properties and
empirical finite-sample performance with that of other commonly used
approximate likelihood approaches, we highlight situations in which the
use of GVA would be advantageous in practice.

\subsection{Estimating abundance in small populations using pedigree
reconstruction}\label{estimating-abundance-in-small-populations-using-pedigree-reconstruction}

\emph{Sarah Croft}\index{Croft, Sarah} \emph{(University of Otago)}

\emph{Co-authors: Jamie Sanderlin, Michael Black, Richard Barker and
Matthew Schofield}

\setlength{\parskip}{0.5em}

Accurate measures of abundance are essential for successful monitoring
of animal populations, and for assessing the efficacy of conservation
interventions. In previous work, genotypic information has been
incorporated into Mark-Recapture models to enable the identification of
individuals, as well as determination of kinship between observed
individuals in Close-Kin Mark-Recapture (CKMR) models. Generally, CKMR
models make large sample assumptions limiting their application to many
endangered and at-risk species. We have developed Bayesian methodology
to estimate population abundance and dynamics for small, isolated
populations using pedigree reconstruction. The true underlying pedigree
completely describes the abundance and population structure over time,
however the true relationships between individuals in wild populations
are rarely known. Given a set of observed genotypes, along with
supplementary data, our methodology is able to successfully reconstruct
the pedigree without the need for large sample assumptions. Prior
knowledge of the mating structure and reproductive dynamics of the
population can also be incorporated in the model. In this talk I will
present our pedigree reconstruction approach for population estimation
using dead recovery data, and will discuss the challenges associated
with the full pedigree approach.

\subsection{Genstat Markdown: Reproducible Research with
Genstat}\label{genstat-markdown-reproducible-research-with-genstat}

\emph{James M. Curran}\index{Curran, James M.} \emph{(University of
Auckland)}

\emph{Co-authors: Yuxiao Wang, Simon Urbanek, and David Baird}

\setlength{\parskip}{0.5em}

Reproducible research ensures that scientific findings can be
independently verified by others using the same data and methods.
Markdown plays a key role in supporting reproducibility by allowing
researchers to combine code, results, and narrative text in a clear,
lightweight format. Tools like R Markdown or Jupyter Notebooks use
Markdown to embed executable code alongside explanations and outputs,
making it easy to document workflows, share analyses, and regenerate
results consistently across different environments.

Genstat is a statistical software package designed for data analysis,
particularly in biometrical research and applications. Many users know
it through its user-friendly interface. However, it also It offers a
powerful scripting language.

In this talk I will describe a collaborative project between students
and staff at the University of Auckland and VSNi Ltd, that allows
Genstat to be used in conjunction with R Markdown to provide flexible
reproducible research capability to the Genstat user community.

\subsection{Multivariate meta-analysis methods for high-dimensional
data}\label{multivariate-meta-analysis-methods-for-high-dimensional-data}

\emph{Alysha M. De Livera}\index{De Livera, Alysha M.} \emph{(La Trobe
University)}

\emph{Co-authors: Luke Prendergast, Robert Staudte, and Jayamini
Liyanage}

\setlength{\parskip}{0.5em}

Multivariate meta-analysis methods for high-dimensional data

Meta-analysis is a statistical method that combines quantitative results
from multiple independent studies on a particular research question or
hypothesis, with the goal of making inference about the population
effect size of interest. Traditional meta-analysis methods have focused
on combining results from multiple independent studies, each of which
has measured an effect size associated with a single outcome of
interest. Modern studies in evidence synthesis, such as those in
biological studies have focused on combining results from studies which
have measured multiple effect sizes associated with multiple correlated
outcomes. We will present a novel, multivariate meta-analysis method for
obtaining summary estimates of the effect sizes of interest for
high-dimensional data, with applications to real and simulated
high-dimensional data. We will discuss advantages, disadvantages and the
statistical challenges, and present a CRAN-based R package for
implementation of the methods.

\subsection{StableMate: a regression framework for selecting stable
predictors across heterogeneous data
environments}\label{stablemate-a-regression-framework-for-selecting-stable-predictors-across-heterogeneous-data-environments}

\emph{Yidi Deng}\index{Deng, Yidi} \emph{(Melbourne Integrative
Genomics; Department of Anatomy and Physiology, The University of
Melbourne; Research School of Finance, Actuarial Studies and Statistics,
The Australian National University)}

\emph{Co-authors: Jiadong Mao, Jarny Choi, Kim-Anh Lê Cao}

\setlength{\parskip}{0.5em}

Inferring reproducible relationships between biological variables
remains a challenge in the statistical analysis of omics data where p
\textgreater{} 10,000 and n \textless{} 500. Methods that identify
statistical associations lack interpretability or reproducibility. We
can address these limitations by inferring stable associations that are
robust to external perturbation on the data. Stable associations can be
an implication in causality since causal relationships are necessarily
stable in some sense (Pearl et al.~2009). Unstable associations can also
be of interests in certain biological applications to study functional
heterogeneity in a biological system.

We developed a new regression framework, StableMate based on the concept
of stabilised regression (SR), which utilise heterogenous data to
enforce stability (Pfister et al.~2021). Given datasets generated from
different environments, such as experiments or disease states,
StableMate 1. identifies stable predictors with consistent functional
dependency with the response across environments. 2. builds a robust
regression model with stable predictors to enforce generalisable
prediction in unseen environments. The ultimate aim is to build
selection ensembles. However, unlike SR that selects stable predictors
by performing stability tests on every possible predictor subset,
StableMate optimizes efficiency with a greedy search based on our
improved stochastic stepwise selection algorithm. In a simulation study,
we show that StableMate outperformed SR for both variable selection and
prediction and significantly reduces running time. In three case studies
of cancer with different omics data types, we show that StableMate can
also address a wide range of biological questions.

\subsection{Fitting integrated species distribution models using
mgcv}\label{fitting-integrated-species-distribution-models-using-mgcv}

\emph{Elliot Dovers}\index{Dovers, Elliot} \emph{(UNSW Sydney)}

\emph{Co-authors: Jakub Stoklosa, David L. Miller, David I. Warton}

\setlength{\parskip}{0.5em}

Integrated species distribution models (ISDMs) are a useful tool for
ecologists, allowing them to use multiple sources of data to infer the
distribution of species across geographic regions. Recent studies have
shown that including latent spatial terms to account for dependence
structures within and between datasets is often crucial to the
performance these models. However, the inclusion of these latent terms
can make the ISDMs technically challenging to fit, and often require
users to learn/adopt bespoke software.

We describe how ISDMs can be fitted using \texttt{mgcv} on \texttt{R} by
using the grouped family feature (\texttt{gfam}). This permits multiple
likelihoods - and hence multiple data types - to be fitted within the
same model by using each dataset to inform the estimation of common
parameters. We additionally show how smoothers over the geographic
coordinates that approximate Gaussian random fields can be used to
incorporate the necessary latent spatial effects. We use
presence/absence and presence-only data to demonstrate that ISDMs fitted
via \texttt{mgcv} have comparable performance to other available
software - through both simulations and in application, where we predict
the occurrence of a tree species in NSW, Australia.

\subsection{Continental-Scale Bayesian Analysis of Acacia Flowering
Phenology: A Novel Framework Integrating Phylogenetic Signal and
Circular
Statistics}\label{continental-scale-bayesian-analysis-of-acacia-flowering-phenology-a-novel-framework-integrating-phylogenetic-signal-and-circular-statistics}

\emph{Owen Forbes}\index{Forbes, Owen} \emph{(CSIRO)}

\emph{Co-authors: Peter H. Thrall, Andrew G. Young}

\setlength{\parskip}{0.5em}

Climate change is driving widespread shifts in plant phenology globally,
with critical implications for ecosystem functioning and species
interactions. While most phenological studies focus on temperate
deciduous species at local scales, Australia's diverse Acacia genus
offers unique opportunities to understand phenological responses to
variations in rainfall and temperature across continental gradients.
Leveraging digitised herbarium collections, we present the largest
phylogenetically-informed analysis of acacia flowering phenology
undertaken to date, spanning 548 species across Australia's full
climatic spectrum.

We developed a novel phenological modelling approach using a Bayesian
hierarchical framework that makes several key innovations. First, we
extend traditional phylogenetic least squares (PGLS) approaches by
incorporating phylogenetic signal as a structured random effect within a
fully Bayesian context, enabling simultaneous estimation of phylogenetic
covariance alongside spatial and temporal dependencies. Second, we
employ von Mises circular distributions to model flowering time,
naturally accommodating seasonal cyclicity while avoiding linear
boundary artifacts. This modular approach also integrates spatial
autocorrelation structure, time-lagged climate covariates, and
phylogenetic relationships within a unified probabilistic framework. Our
analysis of 20,998 herbarium specimens spans 110 years (1910 – 2020)
across Australia, using climate observations from the Australian Gridded
Climate Data, with Köppen climate zones included in the model to
stratify climate-phenology relationships.

This framework successfully quantifies species-specific climate
sensitivities while accounting for phylogenetic constraints. Expected
effect of time-lagged temperature and rainfall effects vary
significantly across Köppen climate zones and species. Species
vulnerability patterns could help identify taxa which could be
potentially at higher risk under projected climate scenarios. This
modular Bayesian approach demonstrates the research potential of natural
history collections for continental-scale ecological inference,
providing a statistical template for global phenological studies
addressing critical climate change impacts.

\subsection{Using a linear mixed model based wavelet transform to model
non-smooth trends arising from designed
experiments}\label{using-a-linear-mixed-model-based-wavelet-transform-to-model-non-smooth-trends-arising-from-designed-experiments}

\emph{Clayton Forknall}\index{Forknall, Clayton} \emph{(Queensland
Department of Primary Industries)}

\emph{Co-authors: Alison Kelly, Yoni Nazarathy and Ari Verbyla}

\setlength{\parskip}{0.5em}

The linear mixed model (LMM) representation of the cubic smoothing
spline is a powerful tool for modelling smooth trends arising from
designed experiments. However, when trends arising from such experiments
are non-smooth, meaning that they are characterised by jagged features,
spikes and/or regions of rapid change approximating discontinuities,
smoothing spline techniques prove ineffective. A technique that has
proven useful for the modelling of non-smooth trends is the wavelet
transform. Existing methods to incorporate the wavelet transform into
the LMM framework are varied, but often share a common limitation; that
is, a reliance on classical wavelet approaches that require observations
to be equidistant and dyadic (\(\log_{2}(n)\) is an integer) in number.
More recently, second generation wavelet methods have been developed,
which overcome the limiting constraints imposed by classical wavelet
approaches, enabling the wavelet transform to be applied to sets of
non-equidistant observations, of any number. We present a method for the
incorporation of these second generation wavelets, namely second
generation B-spline wavelets, into an LMM framework to facilitate the
wavelet transform. Furthermore, using the structure implicit in the
resulting B-spline wavelet basis, we propose extensions to the LMM
framework to enable heterogeneity of the associated wavelet variance
across wavelet scales. This provides a new LMM based method which
enables the flexible modelling of non-smooth trends arising in the
conduct of designed experiments. The proposed method is demonstrated
through application to a data set exhibiting non-smooth characteristics,
that arises from a designed experiment exploring the proteome of barley
malt.

\subsection{Koala Distribution and
Abundance}\label{koala-distribution-and-abundance}

\emph{Scott D. Foster}\index{Foster, Scott D.} \emph{(Data61, CSIRO)}

\emph{Co-authors: Scott D. Foster, Wen-Hsi Yang, Peter Caley, John
McEvoy, David Ramsey, Drew Terasaki Hart, David E. Uribe-Rivera, Romane
H. Cristescu, Eric Vanderdys, Adam McKeown, Corey J. A. Bradshaw, Emma
Lawrence, Andrew Hoskins}

\setlength{\parskip}{0.5em}

The koala (\emph{Phascolarctos cinereus}) is a well-known and studied
Australian marsupial, but the species presents a complex case for
conservation. Currently, most koala conservation efforts focus on
local-scale population estimates, which are often based on expert
opinion or anecdotal evidence, and either precede or ignore most (or
all) available data. In contrast, conservation listing advice and
associated recovery plans require population estimates at the
species-range scale. A data-driven, national-scale population estimate
including all available information will therefore guide effective
management of koala populations by providing high-quality and objective
information. To this end, we have designed nationally consistent survey,
implemented it (with the aid of partners) and analysed the resulting
data (and others). The design uses recently developed techniques
(clustered spatially-balanced designs) whilst the analysis uses emerging
models that incorporate multiple data types (e.g.~point process, binary
and count) that are often-enough collected using different equipment,
protocols and staff. Integrated species distribution models (ISDMs)
have, at their heart, a simple point process but this simplicity still
allows for some complexity in terms of how different types of data can
inform the point process. The model, when fitted to koala data,
indicates that the distribution of koalas is patchy throughout much of
eastern Australia. It also infers that there are more koalas than
previously guessed. Our estimates provide unprecedented evidence to
support nationally consistent and spatially explicit decision-making for
koala conservation, and do so with relevant measures of uncertainty. 

\subsection{A semi-supervised framework for diverse multiple hypothesis
testing
scenarios}\label{a-semi-supervised-framework-for-diverse-multiple-hypothesis-testing-scenarios}

\emph{Jack Freestone}\index{Freestone, Jack} \emph{(Macquarie
University)}

\emph{Co-authors: William S. Noble and Uri Keich}

\setlength{\parskip}{0.5em}

Standard multiple testing procedures are designed to report a list of
discoveries, or suspected false null hypotheses, given the hypotheses'
p-values or test scores. Recently there has been a growing interest in
enhancing such procedures by combining additional information with the
primary p-value or score. In line with this idea, we develop RESET
(REScoring via Estimating and Training), which uses a unique
data-splitting protocol that subsequently allows any semi-supervised
learning approach to factor in the available side information while
maintaining finite sample error rate control. Our practical
implementation, RESET Ensemble, selects from an ensemble of
classification algorithms so that it is compatible with a range of
multiple testing scenarios without the need for the user to select the
appropriate one. We apply RESET to both p-value and competition based
multiple testing problems and show that RESET is (1) power-wise
competitive, (2) fast compared to most tools and (3) able to uniquely
achieve finite sample false discovery rate or false discovery exceedance
control, depending on the user's preference.

\subsection{Fitting Generalised Linear Mixed Models using Sequential
Quadratic
Programming}\label{fitting-generalised-linear-mixed-models-using-sequential-quadratic-programming}

\emph{Peter Green}\index{Green, Peter} \emph{(Tadorna Data Science)}

\setlength{\parskip}{0.5em}

Finding maximum likelihood estimates in generalised linear mixed models
(GLMMs) can be difficult due to the often-intractable integral over the
random effects.

Ensuring convergence can be tricky, especially in binomial GLMMs, and
often multiple optimisers and settings need to be tried to get
satisfactory results. This is exacerbated if you want to use parametric
bootstrap for your inference.

Sequential quadratic programming (SQP) is a method for solving
optimisation problems with non-linear constraints. SQP offers an
alternative option for maximising the Laplace approximation to the GLMM
likelihood, bypassing the need for an inner penalised iterative
reweighted least squares (PIRLS) loop.

This talk discusses the implementation of SQP for GLMMs and compares its
performance to other common approaches.

\subsection{Modelling Species Diversity from Citizen-Science Bird
Counts}\label{modelling-species-diversity-from-citizen-science-bird-counts}

\emph{Graham Hepworth}\index{Hepworth, Graham} \emph{(Statistical
Consulting Centre, The University of Melbourne)}

\emph{Co-authors: Jeremy Silver, William Steele, Al Danger}

\setlength{\parskip}{0.5em}

Bird count data from the Greater Melbourne region over a 25-year period
were recorded by volunteer observers, based on 20-minute surveys. Counts
for a large range of bird species were reported either as counts or
simply as “present”. Our investigation focused on whether there has
been an increase in the presence and abundance of noisy miners, an
endemic species known for its aggressive and territorial behaviour, and
whether this could partly explain apparent declines in other species.

The dataset was augmented with weather, geospatial and ecological
information, and a mixed-effects modelling framework was adopted.
Seasonal effects and long-term trends were accounted for via regression
splines. The influence of different observers and locations were
represented as random effects. GAMMs were fitted for (a) the
presence/absence of noisy miners, (b) the abundance of noisy miners, and
(c) species richness for particular groupings of birds. In modelling
species richness, noisy miner abundance was included as a predictor; we
used multiple imputation to account for the missing data when this
species was reported as “present” but no count was recorded.

There was a small, positive association found between noisy miner
abundance and all-species diversity – a surprising result. A stronger
negative association was found for some species groupings. For many
species groups, diversity tended to decrease with increasing distance
from the nearest native vegetation feature, and with increasing wind
speed. The work provided several statistical and computational
challenges.

\subsection{Data-Adaptive Automatic Threshold Calibration for Stability
Selection}\label{data-adaptive-automatic-threshold-calibration-for-stability-selection}

\emph{Martin Huang}\index{Huang, Martin} \emph{(The University of
Sydney)}

\emph{Co-authors: Samuel Muller, Garth Tarr}

\setlength{\parskip}{0.5em}

Stability selection has gained popularity as a method for enhancing the
performance of variable selection algorithms while controlling false
discovery rates. However, achieving these desirable properties depends
on correctly specifying the stable threshold parameter, which can be
challenging. An arbitrary choice of this parameter can substantially
alter the set of selected variables, as the variables' selection
probabilities are inherently data-dependent. To address this issue, we
propose Exclusion Automatic Threshold Selection (EATS), a data-adaptive
algorithm that streamlines stability selection by automating the
threshold specification process. EATS initially filters out potential
noise variables using an exclusion probability threshold, derived from
applying stability selection to a randomly shuffled version of the
dataset. Following this, EATS selects the stable threshold parameter
using the elbow method, balancing the marginal utility of including
additional variables against the risk of selecting superfluous
variables. We evaluate our approach through an extensive simulation
study, benchmarking across commonly used variable selection algorithms
and static stable threshold values.

\subsection{Model-based assessment of functional and phylogenetic
diversity}\label{model-based-assessment-of-functional-and-phylogenetic-diversity}

\emph{Shaoqian Huang}\index{Huang, Shaoqian} \emph{(School of
Mathematics and Statistics, UNSW)}

\emph{Co-authors: David Warton}

\setlength{\parskip}{0.5em}

\subsection{Topic}\label{topic}

Biodiversity is a quantity of fundamental interest, but measuring it
remains challenging. Two key aspects of diversity are functional
diversity (FD) the extent to which species in a community differ in
their ecological functions, and phylogenetic diversity (PD), the extent
to which species differ in their evolutionary histories.

\subsection{Limitations}\label{limitations}

Commonly used distance-based diversity measures, such as Rao's Q and
Faith's PD, are typically affected by sampling intensity, which refers
to the level of effort invested in collecting samples. Specifically,
these measures confound true diversity changes with changes in species
richness driven by different sampling intensities. In addition, there is
no established statistical framework for analyzing how diversity changes
along environmental gradients.

\subsection{Our Model}\label{our-model}

We propose a model-based assessment of PD and FD changes along
gradients. It is robust to variations in sampling intensity, as the
model explicitly captures the main trend in abundance change, namely,
changes in species richness, and then isolates the residual variation,
the
eta\(-diversity change. We then measure PD or FD change using linear contrasts of the eta\)-diversity
parameters, where the contrasts correspond to the main axes of the
phylogenetic or functional similarity matrix.

\subsection{Simulation Analysis}\label{simulation-analysis}

Based on the presence/absence pattern, we conducted simulations in which
we designed phylogenetic distance matrices with different numbers of
species, modelled on real data. The results show that our approach tends
to outperform methods currently used in ecology.

\subsection{Scalable finite mixture of regression models for clustering
species responses in
ecology}\label{scalable-finite-mixture-of-regression-models-for-clustering-species-responses-in-ecology}

\emph{Francis KC Hui}\index{Hui, Francis KC} \emph{(The Australian
National University)}

\setlength{\parskip}{0.5em}

When modeling species assemblages in ecology, clustering species with
similar responses to the environment can facilitate a more parsimonious
understanding of the assemblage, and improve prediction by borrowing
strength across species within the same cluster. One statistical method
for achieving the above is species archetype models (SAMs), a type of
finite mixture of regression model where species are clustered according
to the \emph{shape} of their environmental response.

In this talk, we introduce approximate and scalable SAMs or asSAMs,
which overcomes some of the current computational drawbacks in fitting
SAMs. We show how asSAMs promotes fast uncertainty quantification via
bootstrapping, along with fast variable selection on archetypal
regression coefficients courtesy of a sparsity-inducing penalty.
Simulation studies and an application to presence-absence records of
over 230 species from the Great Barrier Reef Seabed biodiversity project
demonstrate asSAMs can achieve similar to or better estimation,
selection, and predictive performance than several existing methods in
the literature.

\subsection{Using point cloud data to discover genomic regions
associated with dynamic
height}\label{using-point-cloud-data-to-discover-genomic-regions-associated-with-dynamic-height}

\emph{Colleen H Hunt}\index{Hunt, Colleen H} \emph{(University of
Queensland)}

\emph{Co-authors: David R. Jordan and Emma S. Mace}

\setlength{\parskip}{0.5em}

Plant height in grain sorghum (\emph{Sorghum bicolor L. Moench}),
influences yield, lodging resistance, and harvestability, yet
conventional methods usually measure plant height at maturity, but this
single-time-point measurement misses the growth dynamics across the
season. Tracking and analysing height over time can reveal more about
the genetic factors controlling development rate and pattern. Data from
881 sorghum lines from a diverse population were planted in a partially
replicated design with 1190 plots. UAV-based high-throughput phenotyping
was used to generate weekly plant height measurements from emergence to
flowering. These aerial images were collected and processed using
photogrammetric software to create dense point clouds and canopy height
models for each plot. A logistic regression model within a linear mixed
model framework was applied to effectively depict the sigmoidal growth
pattern typical of sorghum, with parameters indicating maximum height,
growth rate, and inflection point. These biologically relevant
parameters offered a straightforward, quantitative summary of each
genotype’s growth trajectory. The logistic growth parameters were then
used in a genome-wide association study (GWAS) to identify loci linked
to the dynamic aspects of height development, with several significant
associations detected. This is among the first studies to combine
logistic growth modelling of high-throughput UAV data with GWAS to
dissect the genetic control of growth dynamics in sorghum. Our findings
underline the usefulness of logistic regression for modelling crop
growth and showcase the effectiveness of combining UAV phenotyping,
statistical growth modelling, and GWAS to identify genomic regions that
influence developmental trajectories in sorghum.

\subsection{Group Sampling with Imperfect Testing for Biosecurity
Applications}\label{group-sampling-with-imperfect-testing-for-biosecurity-applications}

\emph{Adele Jackson}\index{Jackson, Adele} \emph{(ABARES)}

\emph{Co-authors: Belinda Barnes and Mahdi Parsa}

\setlength{\parskip}{0.5em}

Group sampling, also known as pooled or batch sampling, is a standard
technique in biological sciences and the health sector to use limited
resources efficiently. Common objectives include detecting pest species
and inferring prevalence of disease in communities, livestock or
wildlife. The purpose of this paper is to support the design of robust
group sampling strategies when testing processes are imperfect. We
formulate analytical distributions and statistics for grouped
hypergeometric sampling and its binomial approximation that incorporate
a variety of types of imperfect test. These include tests that respond
to the presence or absence of contaminated material in a group, as well
as PCR and serological testing processes where sensitivity depends on
the number of contaminated items in each group. We also formulate the
Hellinger information of a sampling scheme, which allows us to develop
group sampling design strategies that increase the accuracy of inferred
prevalence. This is an essential component in decision-making during
outbreaks of disease and in operational biosecurity applications. Based
on this work, we estimate leakage through a grouped sampling scheme and
show how accounting for leakage can alter sampling strategies and
improve risk management decisions.

\subsection{The geometry of diet: using projections to quantify the
similarity between sets of dietary
patterns}\label{the-geometry-of-diet-using-projections-to-quantify-the-similarity-between-sets-of-dietary-patterns}

\emph{Beatrix Jones}\index{Jones, Beatrix} \emph{(University of
Auckland)}

\emph{Co-authors: Shiyu Fan, Karen Mumme, John Thompson}

\setlength{\parskip}{0.5em}

Food consumption is complex and high dimensional. Nutrition researchers
measure (or attempt to measure) consumption using food frequency
questionnaires, food recalls, or food records. This generates high
dimensional data which is frequently summarized using principle
components, principle components with rotation, or factor analysis. The
resulting linear combinations are called “dietary patterns.” In this
talk we consider quantifying how similar two different sets of dietary
patterns are, assuming the same set of underlying variables has been
collected. We explore using a multivariate extension of Tucker’s
congruence coefficient for this purpose. Tucker’s congruence
coefficient can be thought of as the length of the projection of one
factor direction onto another, or equivalently the (absolute) cosine of
the angle between them; thus it ranges from 0 to 1. In (eg) two
dimensions, we consider the square root of the area of a unit square
projected from one space onto another; this can be generalized for any
number of dimensions. The measure is symmetric and invariant to
rotation. To contextualize scores on this measure, we compute this
measure of agreement for several datasets from the dietary pattern
validation literature, where the same food questionnaire is given to the
same people a few weeks apart. We also consider the effect of truncating
small loadings, as is common when describing dietary patterns.

\subsection{The performance of Yu and Hoff's confidence intervals for
treatment means in a one-way
layout}\label{the-performance-of-yu-and-hoffs-confidence-intervals-for-treatment-means-in-a-one-way-layout}

\emph{Paul Kabaila}\index{Kabaila, Paul} \emph{(La Trobe University)}

\setlength{\parskip}{0.5em}

Consider a one-way layout and suppose that we have uncertain prior
information that the treatment population means are equal or close to
equal. Yu \& Hoff (2018) extended the ``tail method'' for finding a
confidence interval for a scalar parameter of interest that has (a)
specified coverage probability and (b) relatively small expected length
when this parameter takes values in some given set. They used this
extension to find confidence intervals for these treatment means that
have (a) specified coverage probability individually and (b) relatively
small expected lengths when this uncertain prior information happens to
be correct. They assessed the expected lengths of these confidence
intervals, over the whole parameter space, using a semi-Bayesian
analysis. I describe a revealing alternative assessment of these
expected lengths using a fully frequentist analysis.

Yu, C. \& Hoff, P. (2018) Adaptive multigroup confidence intervals with
coverage. \emph{Biometrika}, 105, 319-335.

\subsection{\texorpdfstring{\texttt{heritable} An R package for
heritability calculations for plant breeding
trials}{heritable An R package for heritability calculations for plant breeding trials}}\label{heritable-an-r-package-for-heritability-calculations-for-plant-breeding-trials}

\emph{Fonti Kar}\index{Kar, Fonti} \emph{(Australian National
University)}

\emph{Co-authors: Emi Tanaka}

\setlength{\parskip}{0.5em}

Understanding how much biological variation is heritable - or passed
down from one generation to next - is central to plant breeding.
Heritability is a useful indictor to assess the genetic gain of
desirable traits such as yield or disease resistance in crop trials. We
created \texttt{heritable}, an R package to streamline the calculation
of heritability from \texttt{asreml} and \texttt{lme4} mixed model
outputs in a user-friendly manner. The package provides up to six
different types of heritability measures and helper functions to compare
between these. Our goal is to support decision makers by providing an
intuitive, open and reproducible workflow in calculating heritability in
R.

\subsection{Visualize your fitted non-linear dimension reduction model
in the high-dimensional data
space}\label{visualize-your-fitted-non-linear-dimension-reduction-model-in-the-high-dimensional-data-space}

\emph{P. G. Jayani Lakshika}\index{Lakshika, P. G. Jayani} \emph{(Monash
University, Australia)}

\emph{Co-authors: Dianne Cook, Paul Harrison, Michael Lydeamore,
Thiyanga S. Talagala}

\setlength{\parskip}{0.5em}

Non-linear dimension reduction (NLDR) techniques such as t-SNE, UMAP,
PHATE, PaCMAP and TriMAP provide a low-dimensional representation of
high-dimensional data by applying a non-linear transformation. The
methods and parameter choices can create wildly different
representations, so much so that it is difficult to decide which is
best, or whether any or all are accurate or misleading. NLDR often
exaggerates random patterns, sometimes due to the samples observed, but
NLDR views have an important role in data analysis because, if done
well, they provide a concise visual (and conceptual) summary of
high-dimensional distributions. To help evaluate the NLDR, we have
developed a way to take the fitted model, as represented by the
positions of points in 2D, and turn it into a high-dimensional wireframe
to overlay on the data, viewing it with a tour. Viewing a model in the
data space is an ideal way to examine the fit. One can see whether it
fits the points everywhere or fits better in some places, or simply
mismatches the pattern. It is used here to help with the difficult
decision on which 2D layout is the best representation of the
high-dimensional distribution, or whether the 2D layout is displaying
mostly random structure. It can also help to see how different layouts
made by different methods are effectively the same summary, or how the
different methods have some particular quirks. This methodology is
available in the R package \texttt{quollr}. We will demonstrate the
technique using single-cell data, particularly to understand cluster
structure.

\subsection{A covariate-adaptive test for replicability across multiple
studies with false discovery rate
control}\label{a-covariate-adaptive-test-for-replicability-across-multiple-studies-with-false-discovery-rate-control}

\emph{Dennis Leung}\index{Leung, Dennis} \emph{(University of
Melbourne)}

\emph{Co-authors: Ninh Tran}

\setlength{\parskip}{0.5em}

Replicability is a gold standard for credible scientific discoveries.
The partial conjunction (PC) p-value proposed by Benjamini and Heller
(Biometrics, 2008), which summarizes individual base p-values obtained
from multiple similar studies, can gauge whether a signal of interest is
replicable. However, when a large set of features are examined by these
studies (such as comparable genomewide association studies conducted by
different labs), testing for their replicated signals simultaneously can
pose a very underpowered problem, due to both the multiplicity
correction required and the inherent limitations of PC p-values. This
power deficiency is particularly severe when replication is demanded for
all studies, the most natural benchmark a practitioner performing
meta-analyses may request.

We propose ParFilter, a procedure that combines the ideas of filtering
and covariate-adaptiveness to power up large-scale testing for
replicated signals as described above. It validly reduces the
multiplicity burden by partitioning the studies into smaller groups and
borrowing between-group information to filter out unpromising features.
Moreover, harnessing side information offered by available covariates,
it trains hypothesis weights to encourage rejections of features more
likely to exhibit replicated signals. We prove its finite-sample false
discovery rate control under standard assumptions on the dependence of
the base p-values across features. In simulations as well as a real case
study on autoimmunity based on RNA-Seq data obtained from thymic cells,
the ParFilter has demonstrated competitive performance against other
existing methods for such replicability analyses.

\subsection{Running Human Subject Experiments via Online
Crowdsourcing}\label{running-human-subject-experiments-via-online-crowdsourcing}

\emph{Patrick Li}\index{Li, Patrick} \emph{(Australian National
University)}

\setlength{\parskip}{0.5em}

Crowdsourcing platforms such as Amazon Mechanical Turk and Prolific
provide scalable and accessible tools for running online human subject
experiments. This talk offers an overview of how to design, set up, and
manage such studies effectively. Drawing on experience from multiple
projects, I will walk through the key steps for obtaining ethics
approval, compare platform workflows, and discuss considerations for
participant recruitment, screening, task design, quality control, cost
estimation, and common pitfalls. This session is intended for
researchers planning to conduct behavioural, perceptual, or
decision-making experiments, as well as those developing data annotation
pipelines for machine learning or applied research.

\subsection{Risk of Guillain-Barré Syndrome after COVID-19 vaccination
and SARS-CoV-2 infection: A multinational self-controlled case series
study}\label{risk-of-guillain-barr-syndrome-after-covid-19-vaccination-and-sars-cov-2-infection-a-multinational-self-controlled-case-series-study}

\emph{Han Lu}\index{Lu, Han} \emph{(University of Auckland)}

\emph{Co-authors: Sharifa Nasreen, Yannan Jiang, Arier Lee, Helen
Petousis-Harris, Jim Buttery, Steve Black, and Jeffrey Kwong}

\setlength{\parskip}{0.5em}

Introduction: The Global Vaccine Data Network (GVDN) is a multinational
collaboration established to conduct globally coordinated
epidemiological studies on vaccine safety using health data. This study
aimed to assess the risk of Guillain-Barré syndrome (GBS) within 42
days after exposure to any COVID-19 vaccines and SARS-CoV-2 infection,
combing multisite data using a common GVDN protocol and data model.

Methods: We used a self-controlled case series (SCCS) design and
identified GBS cases via electronic data sources (EDS) from 20 GVDN
sites globally. Fifteen sites performed medical chart review using
Brighton Collaboration case definitions to determine the level of
diagnostic certainty (LOC). The relative incidence (RI) between
pre-defined risk and control windows was calculated using conditional
Poisson regression, controlling for seasonality. De-identified
case-level data were combined for individual patient data (IPD)
analysis, and the estimates were aggregated with otherwise locally
analysed site-level results to calculate the overall RI with 95\%
confidence interval (95\%CI) using random effects meta-analysis.

Results: We identified 2086 GBS cases (4329 vaccine doses) from EDS and
410 cases were chart-reviewed (12\% with LOC 1-2). We observed an
association between Vaxzevria/Covishield vaccine and GBS among LOC 1-2
cases within 42 days post-exposure (RI 3.10, 95\%CI 1.12-8.62). No
increased risk of GBS was observed with mRNA and other vaccines. Among
489 EDS -identified cases post-infection, significant risk was found
between SARS-CoV-2 and GBS (RI 3.35, 95\%CI 1.83-6.11). The SCCS method
is an efficient approach to collect minimal data for IPD meta-analysis.
Full results using different SCCS models will be discussed.

\subsection{Crossvalidation for predictive models in complex survey
data}\label{crossvalidation-for-predictive-models-in-complex-survey-data}

\emph{Thomas Lumley}\index{Lumley, Thomas} \emph{(University of
Auckland)}

\emph{Co-authors: Amaia Iparragirre}

\setlength{\parskip}{0.5em}

Cross-validation is a standard technique for choosing models with good
out-of-sample prediction error in statistical learning. Chopping data up
at random makes sense for independent observations, but not for
observations from a multistage survey. For example, when a single
cluster is split between test and training sets there is the potential
for data leakage and underestimation of prediction error. I will
describe an approach to cross-validation for complex survey data using
replicate weights and describe its implementation in the R survey
package.

\subsection{Building Trust Without Peer Review: Establishing
Reproducibility Standards in Industrial Statistical
Consulting}\label{building-trust-without-peer-review-establishing-reproducibility-standards-in-industrial-statistical-consulting}

\emph{Dean Marchiori}\index{Marchiori, Dean} \emph{(Wave Data Labs)}

\setlength{\parskip}{0.5em}

In academic research, peer review serves as a cornerstone for ensuring
the credibility and reproducibility of findings. However, in industry,
statistical consultants often operate without the formal structures of
peer review, posing gaps in ensuring the reproducibility and
trustworthiness of their analyses.

This talk explores how statistical consultants for industry can
proactively establish and adhere to standards that promote
reproducibility and foster trust, even in the absence of traditional
peer review mechanisms.

Researchers, data scientists and statisticians often engage with and
provide expertise to government, industry and other groups outside of
academia. This expertise is often trusted implicitly and relied on to
make important decisions. On the other hand, there are many well known
criticisms of conventional peer review systems both in academic
publishing and in commercial work.

Drawing on practices from various industries, we will propose a
pragmatic framework for the development and implementation of
standardized reporting formats, the use of open-source tools for
reproducible analysis, and the adoption of best practices for
documentation and code sharing. We will also examine the role of
professional organizations and industry consortia in setting guidelines
that encourage transparency and accountability.

Attendees will gain insights into practical strategies for implementing
reproducibility standards in their own consulting and research practices
and contribute to a broader movement towards transparency and
accountability in industrial statistics.

\subsection{Accounting for heterogeneous detection rates when inferring
eradication of an invasive
species}\label{accounting-for-heterogeneous-detection-rates-when-inferring-eradication-of-an-invasive-species}

\emph{Sean A. Martin}\index{Martin, Sean A.} \emph{(QUT School of
Mathematics/SAEF)}

\emph{Co-authors: Leah South, Zachary T. Carter, Michael Bode}

\setlength{\parskip}{0.5em}

Eradication of invasive species from islands is an expensive but often
crucial activity for maintaining the ecosystems of these islands. A key
challenge in limiting the costs of an eradication campaign is knowing
the size of the remaining population of invasive animals; this
information informs both ongoing resource planning and when to end a
campaign. Campaign progress is measured by the decline in detections of
the target species. However, detection of a population is imperfect and
highly variable, obscuring the true rate of decline. As a result,
uncertainty arises when detections cease - have all individuals truly
been eradicated, or do some remain undetected? Statistical models have
been developed to infer eradication success based on a record of
detections (sighting records), however, most models ignore natural
heterogeneity in detectability. Our work examines the influence of
variance in detection rates on the false positive and false negative
error rates of hypothetical eradication campaigns and relates this to
overall costs of island conservation. We use a combination of Bayesian
likelihood and likelihood-free models to overcome simultaneous inference
limits on abundance and detectability within this highly stochastic
study system.

\subsection{Spatio-Temporal Species Distribution
Modelling}\label{spatio-temporal-species-distribution-modelling}

\emph{Sam Mason}\index{Mason, Sam} \emph{(UNSW)}

\setlength{\parskip}{0.5em}

Detecting species response to climate change is a critical concern in
ecology requiring relevant climatic predictors over long temporal
windows and large spatial extents. To date this has been technologically
challenging resulting in the widespread use of 30 year averaged datasets
which, while readily accessible, actually mask species responses by
effectively treating the climate as temporally static.

Recent advances in data acquisition technology have made it easier to
obtain a wide range of environmental predictors at appropriate
spatio-temporal scales and at consistent resolutions needed for
effective research.

By using dynamic predictors and anomalies from their mean values in
spatio-temporal species distribution models we see if we can detect
species response to climate change and if this yields better predictive
performance than current practice. Interestingly, we have not found
appreciable spatio-temporal signal, in contrast to expectations and
other work which used naïve approaches to analysis.

\subsection{Multi-environment trial analysis of count data with complex
variance structures using generalised linear mixed
models}\label{multi-environment-trial-analysis-of-count-data-with-complex-variance-structures-using-generalised-linear-mixed-models}

\emph{Michael H. Mumford}\index{Mumford, Michael H.} \emph{(Queensland
Department of Primary Industries)}

\emph{Co-authors: Alan Huang, Alison M. Kelly}

\setlength{\parskip}{0.5em}

The analysis of response data from agricultural field experiments
conducted in multi-environment trials (METs) is typically performed
using linear mixed models (LMMs). The strength of the LMM framework is
in the ability to seamlessly model experimental design terms and to
include complex variance structures, especially for unbalanced data.
When the response variable is count data, the assumptions underpinning
the LMM are violated and it is necessary to extend to a generalised
linear mixed model (GLMM) approach. Statistical modelling using GLMMs
introduces additional complexities arising from a combination of (i)
susceptibility to large estimation biases due to approximations of the
marginal likelihood, (ii) accounting for heterogeneity of
variance/dispersion, and (iii) the increase in computational resources
required. In this talk, a statistical methodology is proposed for the
MET analysis of count data. The analysis approach uses a GLMM framework,
assuming an underlying mean-parameterised Conway-Maxwell Poisson
distribution, that can account for arbitrarily under and over-dispersed
count data. This framework enables partitioning of residual variation
from genetic and other extraneous sources of variation, and adopts a
factor-analytic model for the genotype by environment interaction
effects. The proposed methodology is applied to a series of common bean
trials, where the aim is genotype selection for the response variable
pod count per plant. The analysis is implemented using the
\texttt{glmmTMB\ R}-package, which uses automatic differentiation to
enhance computational speed, the Laplace approximation estimation method
to reduce estimation biases, and a residual maximum likelihood
(REML)-like correction to further reduce estimation biases for variance
components.

\subsection{Better Conversations, Better Support: Strengthening
Consulting through Practical Education and
Community}\label{better-conversations-better-support-strengthening-consulting-through-practical-education-and-community}

\emph{Sharon G. Nielsen}\index{Nielsen, Sharon  G.} \emph{(Adelaide
University)}

\emph{Co-authors: Sam Rogers and Annie Conway}

\setlength{\parskip}{0.5em}

Effective statistical consulting in agriculture depends not just on
technical expertise, but on shared understanding, clear communication,
and the confidence of those seeking advice. Through its series of
practical agronomic workshops and the linked Community of Practice, the
Biometry Hub has shown how targeted education can transform the way
consulting support is delivered and received.

The workshops, typically four core sessions delivered over five days,
cover the design of robust experiments, sound data collection, practical
analysis, and clear interpretation of results. By equipping researchers
and industry staff with these practical skills, the workshops lay the
groundwork for more productive conversations between statisticians and
collaborators.

The Community of Practice then extends this support beyond the workshop
room. Monthly meetings held over Zoom create an informal space for
participants to revisit concepts, tackle real-world questions, and build
confidence alongside peers and mentors. This talk will share how
combining hands-on training, tailored tools like the biometryassist R
package, and an ongoing network of support strengthens the consulting
relationship, helping ensure that good statistical practice is not only
taught, but actively put into practice across the grains industry.

\subsection{Teaching Meta-Analysis for Systematic Reviewers with Mixed
Statistical
Training}\label{teaching-meta-analysis-for-systematic-reviewers-with-mixed-statistical-training}

\emph{Xu Ning}\index{Ning, Xu} \emph{(University of New South Wales)}

\emph{Co-authors: David K E Chan, Coralie Williams, and Peter Humburg}

\setlength{\parskip}{0.5em}

Higher-degree research (HDR) students often consider conducting a
systematic review of their research domain, with the goal of
synthesising their findings with a meta-analysis. However, the
statistical training of HDR students and their supervisors can range
from little to comprehensive. Consequently, we found in our
consultations with these students that they may struggle in specifying
their meta-analytical model, misinterpret their model outputs and
diagnostics, or both, especially with complex meta-analytical models.
Hence, we have developed a short course on meta-analysis to address
these knowledge gaps. The course aims to develop students’ mastery in
meta-analytical methods. We will present empirical findings on whether
the course’s aim was met and discuss the pedagogical considerations
made to accommodate the range of statistical training in our target
cohort.

\subsection{A circular hidden Markov model for directional time series
data}\label{a-circular-hidden-markov-model-for-directional-time-series-data}

\emph{A.A.P.N.M. Perera}\index{Perera, A.A.P.N.M.} \emph{(The Australian
National University)}

\emph{Co-authors: Francis K.C. Hui, Carolyn Huston, A. H. Welsh}

\setlength{\parskip}{0.5em}

Modeling directional time series data such as wind or ocean current
direction presents several interesting challenges. Standard linear time
series techniques do not account for the circularity of the
observations, while existing circular modeling approaches typically work
best when the data span a small arc.

Motivated by a series of fire burn experiments collecting wind direction
in southeastern Australia, we propose a new method for directional time
series data that combines the flexibility of hidden Markov models for
capturing different latent states at different periods of time, with a
conditional von Mises distribution given the latent state to explicitly
account for the circular nature of the responses. The resulting circular
hidden Markov model (cHMM) can allow for multimodality and/or varying
amounts of circular dispersion over time.

Furthermore, by utilizing a von Mises distribution whose mean direction
depends on previous observations, we can accommodate serial correlations
within a specific hidden state. We employ direct maximum likelihood
estimation to fit the cHMM, and examine three approaches to perform
forecasting based on extrapolating the latent state sequence and then
direction observations conditional on this sequence. An application to
the motivating wind direction datasets reveals that cHMMs produce
similar or better point/probabilistic forecasting performance compared
with several established time series methods.

\subsection{Tales from the jungle: a personal perspective of statistical
consulting since
COVID}\label{tales-from-the-jungle-a-personal-perspective-of-statistical-consulting-since-covid}

\emph{Alice Richardson}\index{Richardson, Alice} \emph{(Australian
National University)}

\setlength{\parskip}{0.5em}

The period since COVID-19 has brought unique challenges to statistical
consulting, with researchers venturing into unfamiliar statistical
territory. The activities of an academic statistical consulting practice
have therefore very much taken on the feel of a journey through a
statistical jungle.

In this talk I will describe some of the more alarming encounters I have
experienced in this jungle over the last five years. Examples will be
selected from my client interactions within the broad biometric research
field. Names and topics will be altered to protect the participants, but
not so much as to alter the key learnings from the encounters.

These examples will cover the entire range of statistical endeavour.
From power and sample size calculation through data collection to
modelling, presentation of results and visualisation, no area of
statistical activity is untouched by surprising methodological
suggestions from researchers. The simultaneous emergence of large
language models has also brought unexpected consequences in consulting
practice, as researchers increasingly turn to AI tools for statistical
guidance. I’ll draw out the patterns of common traps and highlight how
the biometric community can assist and advise.

I will also discuss some of the practical strategies I have developed to
carve out a safe path. These strategies can be used for navigating
similar consulting challenges and for fostering better statistical
practices in other collaborative research environments beyond academia.

\subsection{Speed: An R package for Spatially Efficient Experimental
Designs}\label{speed-an-r-package-for-spatially-efficient-experimental-designs}

\emph{Sam Rogers}\index{Rogers, Sam} \emph{(University of Adelaide)}

\emph{Co-authors: Wasin Pipattungsakul, Julian Taylor}

\setlength{\parskip}{0.5em}

Agricultural field trials are typically designed using robust
statistical randomisation, and best practice agricultural field trials
also consider the spatial layouts of the experiment, and how that
interacts with the treatments of interest. However, software providing
access to spatially optimised experimental designs for field trials is
not readily available and can be difficult for new users to get started
with due to sparse documentation.

In this talk, we discuss the newly developed \emph{speed} R package
which provides easy access to a fully open-source package with
comprehensive documentation to enable the design of spatially optimal
experiments. The package offers model-free spatially optimal designs via
a simulated annealing optimisation algorithm. It can produce a spatially
optimal version of many types of experimental designs commonly used in
agricultural research as well as many more complex designs such as
incomplete block designs and partially replicated designs. It provides
multiple objective functions out of the box, with the additional
flexibility to choose or enable custom optimisation metrics, depending
on the objective of the researcher. It also provides some helper
functions for plotting and evaluating experimental designs either
produced via \emph{speed} or alternative design packages.

To demonstrate the package's capabilities, we present spatially optimal
designs for challenging scenarios including two-dimensional blocking and
partially replicated designs. We also show that speed is significantly
faster compared to alternative software.

\subsection{Simulation-based Inference about Variance
Components}\label{simulation-based-inference-about-variance-components}

\emph{Farwa Saleem}\index{Saleem, Farwa} \emph{(School of Mathematics
and Statistics, UNSW Sydney)}

\emph{Co-authors: David Warton}

\setlength{\parskip}{0.5em}

Variance components in generalized linear mixed models (GLMMs) are used
to quantify random effects, for example, they can be used in ecological
studies to quantify the size of the effect of an environmental variable
on species composition. However, constructing accurate confidence
intervals for these components remains challenging with small samples or
near-boundary estimates, where asymptotic methods (e.g., profile
likelihood, Wald intervals) do not perform well. To address this, we
propose a novel simulation-based framework using bootstrap test
inversion coupled with quantile regression. Our approach efficiently
estimates confidence limits by leveraging the continuous distribution of
the estimator, eliminating discretization steps inherent in conventional
bootstrap procedures. This enhances robustness in low-information
scenarios while adapting to boundary constraints. Simulation studies
demonstrate that our method achieves superior coverage accuracy compared
to existing asymptotic and bootstrap alternatives under small-sample
conditions and near-boundary parameters. We show the application of the
technique to an ecological case study analysing response of species
composition to environmental gradients. This work delivers a flexible
tool for more defensible inference in applications demanding precise
uncertainty quantification for random effects.

Keywords: Variance components, GLMM, bootstrap, quantile regression,
confidence intervals, ecological statistics.

\subsection{The 4S method for the longitudinal analysis of
multidimensional questionnaires: application to Parkinson’s disease
progression from patient
perception}\label{the-4s-method-for-the-longitudinal-analysis-of-multidimensional-questionnaires-application-to-parkinsonuxe4uxf4s-disease-progression-from-patient-perception}

\emph{Tiphaine Saulnier}\index{Saulnier, Tiphaine} \emph{(New Zealand
Brain Research Institute, Christchurch, New Zealand and Institute of
Neurodegenerative Diseases (IMN), University of Bordeaux, Bordeaux,
France)}

\emph{Co-authors: Cécile Proust-Lima, Daniel J. Myall, Toni L. Pitcher,
Leslie Livingston, Rachel Nolan, Sophie Grenfell, John C.
Dalrymple-Alford, Alexandra Foubert-Samier, Wassilios G. Meissner, Tim
J. Anderson and Campbell J. Le Heron}

\setlength{\parskip}{0.5em}

\textbf{Introduction} In health research, questionnaires are widely used
to assess clinical states or quality of life (QoL). Longitudinal
analysis of these tools offers valuable insights into disease
progression. However, this raises multiple challenges, such as handling
repeated ordinal responses, capturing the multidimensional traits
underlying all the items, and accounting for informative dropout due to
death during follow-up. In this work, we present the 4S method – a
comprehensive strategy developed to address these challenges, and
illustrate it describing health-related QoL (Hr-QoL) changes in
Parkinson’s disease (PD).

\textbf{Methods} The 4S method comprises four successive steps: (\emph{1
– structuring}) identify questionnaire dimensions through factor
analyses; (\emph{2 – sequencing}) describe each dimension progression
and associated predictors with a joint latent process model combining an
item-response mixed model and a proportional hazard for death risk;
(\emph{3 – staging}) compare progression across dimension continuums
by projecting clinical stages; (\emph{4 – selecting}) highlight
stage-specific most informative items based on Fisher information.

\textbf{Application} We analyzed longitudinal data from the New Zealand
Parkinson Progression Programme (NZP\(^3\)), following over 400 PD
patients for up to 16 years. Hr-QoL was measured via the PDQ-39
questionnaire, covering motor and non-motor spheres. Five dimensions
were identified: mobility, daily activities, psycho-social, stigma, and
cognition/communication/bodily discomfort. All, except stigma, showed
progressive decline, with patterns notably varying by sex and onset age.
Items related to walking, dexterity, anxiety, and communication appeared
as particularly sensitive during PD stages, offering guidance to enhance
patient-centered care.

\textbf{Conclusion} The 4S method is a comprehensive statistical
strategy suited to analyze repeatedly-collected questionnaires in health
studies.

\subsection{Enhancing Fraud Detection in Banking through Random Survival
Forests: Addressing Data Imbalance and Model
Transparency}\label{enhancing-fraud-detection-in-banking-through-random-survival-forests-addressing-data-imbalance-and-model-transparency}

\emph{Arjun Sekhar}\index{Sekhar, Arjun} \emph{(Macquarie University)}

\emph{Co-authors: Samuel Muller and Maurizio Manuguerra}

\setlength{\parskip}{0.5em}

With the rise in financial fraud amid growing transaction volumes, in
this presentation we share our insights from the use of Random Survival
Forests (RSFs) for enhancing fraud detection in banking systems. While
many machine learning techniques prioritise accuracy, they overlook two
crucial issues: the imbalance of fraud versus legitimate transaction
data, and the transparency of the model decision-making process.

We position RSFs as a potential solution, leveraging their time-to-event
structure to enable dynamic fraud prediction and their ensemble nature
to handle high-dimensional, skewed datasets. Using real-world
transactional data, the presentation addresses model imbalance through
resampling strategies and evaluates the model’s ability to detect
minority-class fraud events without inflating false positives.

Furthermore, the interpretability of RSFs is explored through
explainable AI frameworks, offering transparency essential for
regulatory compliance and institutional trust. By investigating RSFs to
deliver fraud detection that is both accurate and accountable, this
presentation bridges algorithmic robustness with transparency, advancing
a practical and governance-aligned solution for responsible AI
deployment in today’s increasingly regulated financial services
landscape.

\subsection{Variable Selection in a Joint Model for Huntington's Disease
Data}\label{variable-selection-in-a-joint-model-for-huntingtons-disease-data}

\emph{Rajan Shankar}\index{Shankar, Rajan} \emph{(The University of
Sydney)}

\emph{Co-authors: Tanya Garcia and Garth Tarr}

\setlength{\parskip}{0.5em}

Huntington's disease is a neurodegenerative disease caused by a
defective Huntingtin gene, with symptoms that progressively worsen and
eventually lead to a clinical diagnosis. Identifying the clinical and
demographic factors that influence symptom severity and
time-to-diagnosis is critical for understanding disease progression so
that early-intervention strategies can be timely implemented. We propose
a joint model to relate symptom severity \(y\) and time-to-diagnosis
\(x\), conditional on clinical and demographic predictor variables
\(\mathbf{z}\). However, it may be that certain predictor variables are
important for \(y\) but not for \(x\) and vice-versa, so we use
regularisation techniques to select different sets of predictor
variables for \(y\) and \(x\). Since \(x\) is a time-to-event variable,
there is the added challenge that many of its values are right-censored
due to individuals who did not develop the disease during the study.
Therefore, to fit the joint model, we apply the expectation-maximisation
(EM) algorithm to alternate between parameter estimation and imputation
of the right-censored values until convergence. We demonstrate our
method on Huntington's disease patient data, showcasing how users can
choose appropriate values for the regularisation tuning parameters.

\subsection{An allometric differential equation model quantifies energy
trade-offs between growth and reproduction under temperature
variation}\label{an-allometric-differential-equation-model-quantifies-energy-trade-offs-between-growth-and-reproduction-under-temperature-variation}

\emph{Hideyasu Shimadzu}\index{Shimadzu, Hideyasu} \emph{(Kitasato
University)}

\setlength{\parskip}{0.5em}

Individual variation in growth and reproduction is common even within a
single species. Such differences play a crucial role in shaping
life-history strategies, particularly in response to environmental
factors such as temperature. This talk introduces a novel system of
allometric differential equations that models internal energy allocation
between two fundamental biological processes, growth and reproduction,
under different thermal regimes.

Climate change is not only driving increases in mean temperatures but
also intensifying temperature variability, resulting in greater
environmental unpredictability. While temperature is known to regulate
resource allocation in ectotherms, the consequences of stochastic
thermal fluctuations for life-history traits remain poorly understood.
Using \emph{Daphnia magna} as a model organism, we present a novel
allometric growth model that characterises energy allocation dynamics in
fluctuating thermal environments. Our framework quantifies the effects
of randomly varying temperatures on lifetime patterns of growth and
reproduction. Our model reveals that exposure to unpredictable thermal
regimes can elicit life-history responses similar to those observed
under persistently elevated temperatures. Importantly, our energy-based
approach identifies patterns of reproductive investment that are not
discernible from growth data alone, highlighting the pivotal role of
temperature variability in shaping life-history trajectories. As climate
change increasingly entails unpredictable environmental conditions, the
model and findings presented here offer valuable tools for anticipating
and managing biological responses to future climate scenarios.

\subsection{Zero-inflated Tweedie distribution for abundance of rare
ecological
species}\label{zero-inflated-tweedie-distribution-for-abundance-of-rare-ecological-species}

\emph{Nokuthaba Sibanda}\index{Sibanda, Nokuthaba} \emph{(Victoria
University of Wellington)}

\setlength{\parskip}{0.5em}

Abundance data for rare species can be extremely zero-inflated, where
percentage of zeros can be over 90\%. This poses a challenge even for
the standard Tweedie model which naturally allows for a probability mass
at zero with continuous non-negative values. We investigate use of a
zero-inflated Tweedie distribution when modelling non-negative
continuous abundance values with zeros for rare species. Despite their
significance, research on zero-inflated models has predominantly focused
on count models such as zero-inflated Poisson or negative binomial
regressions, with only recent studies exploring zero-inflated Tweedie
models in insurance claims (Zhou, Qian, and Yang 2022; Gu2024; So and
Valdez 2024; So and Deng 2025).

The zero-inflated Tweedie model uses a mixture model approach that
integrates a Tweedie model with a binary model to distinguish between
excess zeros - those resulting from an independent process, and true
zeros those resulting from the Tweedie model itself. We use a Bayesian
approach to estimate the model parameters. We model means of the Tweedie
model using a log-link function with covariates and unobserved random
effects. The spatial association between observations is accounted for
using a conditionally auto-regressive prior.

\subsection{A missing data detective story – how I navigated through a
perfect storm of drop-out’s, COVID and informative
missingness.}\label{a-missing-data-detective-story-uxe4uxec-how-i-navigated-through-a-perfect-storm-of-drop-outuxe4uxf4s-covid-and-informative-missingness.}

\emph{Eve Slavich}\index{Slavich, Eve} \emph{(University of New South
Wales)}

\emph{Co-authors: Rebecca Young, Anita Pfeiffer, Emmanuel Karantanis}

\setlength{\parskip}{0.5em}

Real-world clinical studies often present missing data challenges that
leave the textbook educated scratching their heads. Through consulting,
I become involved in \href{https://doi.org/10.1016/j.cont.2025.101762}{a
study} where multiple correlated missing data mechanisms created
analytical challenges. Following implementation of new birthing
protocols aimed at reducing birth-related injuries, in a Sydney health
district, a follow-up study was designed to evaluate incontinence
outcomes. However, the missing data pattern proved to be informatively
related to both treatment and outcome variables: women who experienced
successful outcomes (no incontinence) were systematically less likely to
attend follow-up appointments. The situation was further complicated
when the COVID-19 pandemic struck mid-study, dramatically reducing
attendance rates and introducing an additional layer of correlated
missingness that varied by treatment timing and baseline
characteristics. Working within the constraints of the study design and
the budget for a statistician, I developed a model for missingness which
I will outline in this talk.

While missing data complications can severely challenge study validity,
thoughtful statistical approaches can still yield meaningful insights
when the caveats (caveats galore!) are properly acknowledged and
addressed.

References: Young, R et al, 2025, Outcomes for obstetric anal sphincter
injuries and anal incontinence following introduction of a perineal
bundle, Continence, Volume 14

Keywords: Missing data analysis, informative missingness, multiple
imputation, statistical consulting, statistical practice

\subsection{Rate-optimal sparse gamma scale mixture
detection}\label{rate-optimal-sparse-gamma-scale-mixture-detection}

\emph{Michael Stewart}\index{Stewart, Michael} \emph{(The University of
Sydney)}

\emph{Co-authors: Qikun Chen}

\setlength{\parskip}{0.5em}

We consider a model where observations from a known gamma distribution
are possibly contaminated by observations from another gamma
distribution with the same shape but a different mean. Such a model has
been considered for times between neurotransmitter releases based on a
Markov chain with amalgamated indistinguishable states. We focus on the
case where the contaminating component occurs rarely, the so-called
sparse gamma scale mixture detection problem. Due to the irregularity of
such models theoretical results concerning detectability bounds are
non-standard. Nonetheless in recent years a body of theory has been
developed which covers the case when the mean of the unknown
contaminating component is smaller than the null mean, but not when it
is larger. We present some recent results filling this gap in the
literature. In particular we describe a test which attains the optimal
rate of convergence in various local alternative scenarios which is a
Bonferroni-type test combining three different tests.

\subsection{False Discovery Rate Controlled Robust Variable Selection
under Cellwise
Contamination}\label{false-discovery-rate-controlled-robust-variable-selection-under-cellwise-contamination}

\emph{Xiaoya Sun}\index{Sun, Xiaoya} \emph{(School of Mathematical and
Physical Sciences, Macquarie University)}

\emph{Co-authors: Houying Zhu, Garth Tarr and Samuel Muller}

\setlength{\parskip}{0.5em}

The increasing complexity and dimensionality of modern datasets make
cellwise outliers a frequent and expected phenomenon, complicating tasks
such as variable selection. An mRNA-Seq expression profiling data set of
human brain tissue for Huntington's Disease (HD) illustrates this
practically, where we work to identify key genes involved in HD by
comparing affected individuals with neurologically normal controls. Our
objective is to obtain reliable results while minimizing the impact of
cellwise outliers. While many robust variable selection methods have
been developed, few effectively handle high-rate cellwise contamination,
and it remains an open problem to what extent such methods can control
error rates. We introduce GRALF, the Gaussian-Ranked Adaptive Lasso with
False discovery rate (FDR) control, a robust method designed for
high-dimensional variable selection under cellwise contamination, taking
advantage of FDR control to increase statistical reliability. GRALF
builds on the framework of GR-ALasso, an effective robust variable
selection method using the Gaussian-Rank estimator in the Adaptive
Lasso, and integrates FDR control by generating fake counterparts of the
Gaussian-rank transformed variables and estimating the number of false
discoveries in the optimization process. Simulation studies demonstrate
GRALF's desirable performance in terms of empirical power and FDR
control under various conditions. The analysis of the HD data set
identified potential gene markers associated with the development of HD,
which overlap with findings from previous research, further validating
the effectiveness of GRALF.

\subsection{A Test for Detecting Multiple Clusters with Hotspot Spatial
Properties}\label{a-test-for-detecting-multiple-clusters-with-hotspot-spatial-properties}

\emph{Kunihiko Takahashi}\index{Takahashi, Kunihiko} \emph{(Institute of
Science Tokyo)}

\emph{Co-authors: Hideyasu Shimadzu}

\setlength{\parskip}{0.5em}

Various statistical tests have been widely used in spatial epidemiology
to investigate regional patterns in disease occurrence, particularly to
assess whether disease risk is significantly elevated in specific areas
compared to neighboring regions or adjacent time periods. One such
method is the cluster detection test (CDT), which identifies non-random
spatial distributions of diseases and highlights high-prevalence regions
without prior assumptions. Among CDT methods, scan statistics are
compelling and use a maximum likelihood framework to search across
spatial and/or temporal windows for potential clusters. Examples include
Kulldorff’s circular scan statistic and the flexibly shaped scan
statistic by Tango and Takahashi. More recently, Takahashi and Shimadzu
proposed a scan-based method that simultaneously detects multiple
clusters by integrating generalized linear models and an information
criterion to determine the optimal number of clusters. Traditional
scan-based tests often assume that disease risk is uniformly elevated
within a single cluster. However, they may mistakenly combine multiple
adjacent hotspots—each with potentially different risk levels—into
one, thereby masking meaningful spatial heterogeneity. In this study, we
propose a new test procedure that more accurately identifies adjacent
hotspot clusters as distinct entities. Our approach enhances the scan
statistic framework by incorporating Cochran’s Q-statistic to assess
heterogeneity within clusters. We demonstrate the effectiveness of the
proposed method through real-world applications and compare its
performance with conventional scan-based tests.

\subsection{Integrating Spatial Data and On-Farm Experimentation to
Understand Wheat Variety Performance Across Western
Australia}\label{integrating-spatial-data-and-on-farm-experimentation-to-understand-wheat-variety-performance-across-western-australia}

\emph{Sandra K. Tanz}\index{Tanz, Sandra K.} \emph{(Centre for Crop and
Disease Management, Curtin University)}

\setlength{\parskip}{0.5em}

Understanding spatial variability in crop performance is critical to
improving the statistical design and interpretation of variety trials at
the farm scale. This pilot study investigates the integration of
large-scale on-farm experimentation (OFE) with spatial data analysis to
better understand genotype-by-environment interactions in wheat. The
trial spans five grower-managed sites across Western Australia,
incorporating two core wheat varieties—Scepter (a widely grown
commercial line) and IGW6993 (a near-release InterGrain line)—with
some growers also including Rockstar and Tomahawk CL Plus based on local
relevance.

Each trial is embedded within commercial paddocks and implemented using
strip trial designs aligned with paddock production zones identified
from historical yield maps. This design allows statistical comparisons
of varietal performance across contrasting spatial zones defined by
environmental variation in soil, topography, and management history. The
study aims to (1) develop robust methods for spatially aware
experimental design in commercial cropping systems, and (2) assess the
potential for OFE to complement small-plot trials in supporting variety
selection and agronomic decision-making.

While the trial was only recently seeded (May 2025), preliminary outputs
include trial design validation and early NDVI imagery from emergence.
Future analyses will incorporate additional covariates such as soil
electrical conductivity (EM mapping), NDVI time series, and topographic
indices to model spatial responses. This work highlights the value of
integrating spatial data into trial workflows and explores the role of
predictive breeding tools in commercial agriculture.

\subsection{Outlier-robust estimation of state-space models using a
penalised
approach}\label{outlier-robust-estimation-of-state-space-models-using-a-penalised-approach}

\emph{Garth Tarr}\index{Tarr, Garth} \emph{(The University of Sydney)}

\emph{Co-authors: Rajan Shankar, Ines Wilms and Jakob Raymaekers}

\setlength{\parskip}{0.5em}

State-space models are a broad class of statistical models for
time-varying data. The Gaussian distributional assumption on the
disturbances in the model leads to poor parameter estimates in the
presence of additive outliers. Whilst there are ways to mitigate the
influence of outliers via traditional robust estimation methods such as
M-estimation, this issue is approached from a more modern perspective
that utilises penalisation. A shift parameter is introduced at each
timepoint, with the goal being that outliers receive a non-zero shift
parameter while clean timepoints receive a zero shift parameter after
estimation. The vector of shift parameters is penalised to ensure that
not all shift parameters are trivially non-zero. Apart from making it
feasible to fit accurate and reliable time series models in the presence
of additive outliers, other benefits of this approach include automatic
outlier flagging and visual diagnostics to provide researchers and
practitioners with better insights into the outlier structure of their
data. We will demonstrate the utility of this method on animal tracking
data.

\subsection{A Proportional Random Effect Block Bootstrap for Highly
Unbalanced Clustered
Data}\label{a-proportional-random-effect-block-bootstrap-for-highly-unbalanced-clustered-data}

\emph{Zhi Yang Tho}\index{Tho, Zhi Yang} \emph{(Australian National
University)}

\emph{Co-authors: Alan H. Welsh and Raymond Chambers}

\setlength{\parskip}{0.5em}

Clustered data arise naturally in a wide range of scientific and applied
research settings where units are grouped within cluster, and are
commonly analyzed using linear mixed model to account for within-cluster
correlations. This article focuses on the scenario in which cluster
sizes might be highly unbalanced, and proposes a proportional random
effect block bootstrap that is applicable in such case and is robust to
misspecification of the stochastic assumptions of the linear mixed
model. The proposed method is a generalization of the random effect
block bootstrap, originally designed for the balanced case, and can be
used to perform inferences on parameters of the linear mixed model or
functions thereof. We establish asymptotic consistency of the proposed
bootstrap under general cluster sizes scenario, showing that the
original random effect block bootstrap is only consistent when cluster
sizes are balanced. A modified random effect block bootstrap is also
proposed which enjoys similar asymptotic consistency properties as the
proportional random effect block bootstrap. Simulation study
demonstrates the strong finite sample inferential performance of the
proposed bootstraps, particularly compared with the random effect block
bootstrap and several existing bootstrap methods for clustered data. We
apply the proposed bootstraps to the Oman rainfall enhancement trial
dataset with cluster sizes ranging from 1 to 58. Results show that the
bootstrap confidence intervals based on our proposed bootstraps are more
adequate than those of random effect block bootstrap and that the
employed ionization technology has a statistically significant effect in
increasing the amount of rainfall.

\subsection{Functional Data Analysis for the Australian Grains
Industry}\label{functional-data-analysis-for-the-australian-grains-industry}

\emph{Braden J. Thorne}\index{Thorne, Braden J.} \emph{(Curtin
University)}

\emph{Co-authors: Adam H. Sparks}

\setlength{\parskip}{0.5em}

When performing statistical analysis on time series data with regular
sampling, it is common to apply filtering or windowing methods to reduce
the complexity of the task. While this process enables many classical
approaches, it inevitably leads to compression of the full information
available. An alternative approach is to treat observations as samples
of a continuous mathematical function and focus analysis on the curves
these functions produce rather than the samples. This is the underlying
idea of functional data analysis, a statistical analysis approach that
has seen growing attention in recent years. In this talk I will offer an
introduction to functional data analysis and detail our exploration of
these methods for application to the grains industry across Australia.
Specifically, I will present two case studies; estimating charcoal rot
prevalence in sixteen years of data from in-paddock soybean experimental
trials using weather data, and analysing frost risk in broadacre crops
with varying stubble management practices using sensor data.

\subsection{Extending diagnostic validity meta-analysis to several
diagnostic
guidelines}\label{extending-diagnostic-validity-meta-analysis-to-several-diagnostic-guidelines}

\emph{Alain C. Vandal}\index{Vandal, Alain C.} \emph{(The University of
Auckland/Waipapa Taumata Rau)}

\emph{Co-authors: Henry C. K. Kwok, Fransiska R. M. Falconer, Andrew G.
Hill, Andrew D. MacCormick}

\setlength{\parskip}{0.5em}

Optimal management of acute biliary disease should include an assessment
for possible choledocholithiasis (stones in the bile duct). Various
diagnostic guidelines have been developed by expert bodies for this
purpose, but uncertainties remain about their performance in wider
practice. We carried out a meta-analysis to compare three of these
guidelines. Each guideline yields a score, for which different
thresholds are considered in terms of sensitivity and specificity. Of
1,892 records, thirty-one studies were identified for this
meta-analysis. All studies focused on one or more of three international
guidelines, namely the ASGE guidelines from 2010, the revised ASGE
guidelines from 2019 and the ESGE guidelines from 2019. Each study
reported on the diagnostic validity of one or more of these guidelines
at one or more score cutpoint. The analysis was carried out using the
\emph{diagmeta} R package {[}Rücker G, Steinhauser S, Kolampally S,
Schwarzer G (2022){]}. Their approach allows the meta-analytical
estimation of the area under the receiver operating curve, a quantity
that summarises the performance of a diagnostic tool. The approach
consists in fitting transformed values of sensitivity and specificity at
various cutpoints, which can differ between studies, accounting for
study heterogeneity using linear mixed effects modelling. While the
\emph{diagmeta} approach applies to the meta-analysis of a single
diagnostic procedure, we show how its framework can be extended to
jointly and severally test for performance similarity between
guidelines. We also show how it can be extended to identify outlier
studies, for the purpose of carrying out sensitivity analyses.

\subsection{Integrated Species Distribution Models: A Single-Index
Approach}\label{integrated-species-distribution-models-a-single-index-approach}

\emph{Quan Vu}\index{Vu, Quan} \emph{(Australian National University)}

\emph{Co-authors: Francis Hui, Alan Welsh, Samuel Muller, Eva Cantoni
and Chris Haak}

\setlength{\parskip}{0.5em}

In ecology and fisheries, species abundance data are often collected
using a number of surveys that have different characteristics. Spatial
statistical models can be used to integrate information from these
datasets to improve both interpretability and prediction of the
abundance. In this talk, we introduce the single-index integrated
species distribution model, based on a single index representing the
latent spatial distribution of the species abundance, and
survey-specific link functions representing the different catchability
properties of each survey. We demonstrate the use of the model through
an analysis of scallop data collected from two surveys: a bottom trawl
survey which covers a wide spatial domain but is less efficient, and a
hydraulic dredge survey which is more efficient and spatially targeted.
Results show that our model offers meaningful interpretations of
covariate effects, spatial distribution, and survey catchability
differences, and achieves superior predictive performance compared to
contemporary species distribution models.

\subsection{Extension of the corrected score estimator in a Poisson
regression model with a measurement
error}\label{extension-of-the-corrected-score-estimator-in-a-poisson-regression-model-with-a-measurement-error}

\emph{Kentarou Wada}\index{Wada, Kentarou} \emph{(Tokyo University of
Science)}

\emph{Co-authors: T. Kurosawa}

\setlength{\parskip}{0.5em}

Kukush et al.~(2004) discussed the bias of the naive estimator for the
regression parameters in a Poisson regression model with a measurement
error for the case where the explanatory variable and measurement error
follow normal distributions. Wada and Kurosawa (2023) proposed the
corrected naive (CN) estimator as a consistent estimator for a Poisson
regression model with a measurement error for the case where the
explanatory variable and measurement error are general distributions.
The CN estimator directly calibrates the bias of the naive estimator.
The CN estimator is given by the solution of the estimation equation of
the Poisson regression model under the error-in-variables framework.
However, the CN estimator does not always have an explicit expression
under the condition that the explanatory variable and measurement error
follow general distributions. On the other hand, Kukush et al.~(2004)
considered the corrected score (CS) estimator as a consistent estimator
for the true parameter of the Poisson regression model with a
measurement error. In this research, we extend the CS estimator to the
case where the explanatory variable and measurement error are general
distributions. The new estimator can be applied for the condition that
the CN estimator does not have an explicit expression. As illustrative
examples, we give simulation studies to verify the effectiveness of the
new estimator.

\subsection{Simultaneous Inference for Latent Variable Predictions in
Factor Analytic
Models}\label{simultaneous-inference-for-latent-variable-predictions-in-factor-analytic-models}

\emph{Zhining Wang}\index{Wang, Zhining} \emph{(Research School of
Finance, Actuarial Studies and Statistics, Australian National
University)}

\emph{Co-authors: Emi Tanaka, Francis K.C. Hui}

\setlength{\parskip}{0.5em}

Factor analytic models (also known as latent variable models) are
fundamental tools in multivariate statistics, widely applied in fields
such as psychology, economics, and social sciences. While considerable
attention has been given to estimation and inference on parameters such
as the loading matrix and the error variances, relatively less research
has been done on how to perform inference on the predicted factors,
e.g., how to construct prediction intervals for the latent variables
jointly across all the clusters in a given dataset. We explore a
framework for the simultaneous inference of the predicted latent
variables in factor analytic models. We demonstrate the construction of
simultaneous prediction intervals for the predicted factors, examining
strategies such as the bootstrap and Monte Carlo simulation, and show
how this also facilitates joint/multiple testing across different
cluster-level predictions. We examine the practical feasibility and
robustness of the proposed simultaneous inference methods with
simulation studies and an application in the biosciences.

\subsection{Estimating extinction time from the fossil record using
regression
inversion}\label{estimating-extinction-time-from-the-fossil-record-using-regression-inversion}

\emph{David I. Warton}\index{Warton, David I.} \emph{(UNSW Sydney)}

\emph{Co-authors: Victor Tsang}

\setlength{\parskip}{0.5em}

An important problem in palaeoecology is estimating the extinction or
invasion time of a species from the fossil record - whether because this
is of interest in and of itself, or in order to understand the causes of
extinctions and invasions, for which we need to know when they actually
happened. There are two main sources of error to contend with - sampling
error (because the last time you see a species need not be the last time
it was there) and measurement error (dating specimens, usually well
known). The paleobiology literature typically ignores one or other of
these sources of error, leading to bias and underestimation of
uncertainty to an extent that is often qualitatively important.

The problem is surprisingly difficult to address statistically, because
while standard regularity conditions are technically satisfied, we are
typically close to a boundary where they break down, and hence standard
asymptotic approaches to inference typically perform poorly in practice.
We propose using a novel method, which we call regression inversion, for
exact inference, and we apply this technique to a compound
uniform-truncated t (CUTT) model for fossil data. We show via simulation
that this approach leads to unbiased estimators, and accurate interval
inference, in contrast to its competitors. We show how to check the CUTT
assumption visually, and provide software to apply all of the above in
the reginv package.

\subsection{\texorpdfstring{The \emph{equalto} covariance structure for
meta-analysis using the \texttt{glmmTMB} R
package}{The equalto covariance structure for meta-analysis using the glmmTMB R package}}\label{the-equalto-covariance-structure-for-meta-analysis-using-the-glmmtmb-r-package}

\emph{Coralie C. Williams}\index{Williams, Coralie C.} \emph{(University
of New South Wales)}

\emph{Co-authors: Maeve McGillycuddy, David I. Warton, Shinichi
Nakagawa}

\setlength{\parskip}{0.5em}

Meta-analysis is a widely used statistical method for synthesising
quantitative results from related studies, enabling researchers to
address broad questions and explore sources of heterogeneity. In R
statistical software, several packages support meta-analytic modelling,
including the actively maintained \texttt{metafor} package which is
commonly used across all disciplines. But meta-analysis typically uses a
mixed model, with the variance structure of study-level random effects
known, so it should be possible to fit such models using standard mixed
modelling software. We propose using the \texttt{glmmTMB} package for
this purpose. In the last decade, \texttt{glmmTMB} has developed into a
flexible package to fit generalised linear mixed models (GLMMs) via the
Template Model Builder (TMB) framework. The package holds a similar
interface as the \texttt{lme4} package but supports a broader range of
distributions and covariance structures using Laplace approximation.
Here, we introduce a new covariance structure, called
“\emph{equalto}”, to the \texttt{glmmTMB} package, so that it can be
used to fit meta-analytic models by allowing users to specify a known
variance-covariance matrix for the sampling error. This enables explicit
modelling of heteroscedasticity or dependence among effect sizes. The
new implementation offers an alternative way to do a meta-analysis,
convenient for users already familiar with fitting mixed models in
\texttt{lme4} or \texttt{glmmTMB} packages, and capable of handling
complex model structures. This novel implementation supports more
flexible modelling of meta-analytical data, expanding the R toolkit
available for evidence synthesis. We showcase its applicability with
illustrative examples in ecology and evolution.

\subsection{Disease cluster detection via functional additive models
incorporating spatial
correlation}\label{disease-cluster-detection-via-functional-additive-models-incorporating-spatial-correlation}

\emph{Michio Yamamoto}\index{Yamamoto, Michio} \emph{(The University of
Osaka / RIKEN AIP / Shiga University)}

\emph{Co-authors: Tatsuhiko Anzai, Kunihiko Takahashi}

\setlength{\parskip}{0.5em}

Detecting spatial clusters of diseases is crucial for understanding
disease patterns and developing effective prevention and treatment
strategies. Spatial scan statistics are powerful tools for detecting
spatial clusters with a variable scanning window size. If covariates are
related to an outcome and not geographically randomly distributed,
searching for spatial clusters may require adjusting for the covariates.
In addition, spatial correlation in the outcome, which is often
overlooked during cluster detection, can affect the results. In this
study, we propose a new spatial scan statistic that handles multiple
functional covariates indicating past information over time and the
spatial correlation of the outcome. Our method flexibly models these
factors in the framework of functional additive models. We develop an
optimization algorithm to estimate the model parameters for the normal
outcome case. A simulation study and real data analysis indicate that
the proposed method can detect disease clusters despite longitudinal
covariates and spatial correlations compared to existing methods.

\subsection{Childhood Risk and Resilience Factors for Pasifika Youth
Respiratory Health: Accounting for Attrition and
Missingness}\label{childhood-risk-and-resilience-factors-for-pasifika-youth-respiratory-health-accounting-for-attrition-and-missingness}

\emph{Siwei Zhai}\index{Zhai, Siwei} \emph{(University of Auckland)}

\setlength{\parskip}{0.5em}

In New Zealand, 7\% of deaths are related to respiratory diseases, with
Pacific people at higher risk. Our work investigates the causal effects
of early-life risks and resilience factors on early-adulthood lung
function amongst Pacific Islands Families Study (PIFS) cohort members
(n=1,398). 466 from the cohort participated in the respiratory study.
Primary outcome was forced expiratory volume in 1 second (FEV1) z-score
at age 18 years. FEV1 and healthy lung function (HLF), defined as the
z-score being larger than -1.64, were secondary outcomes. A previous
study had evaluated the effects of early-life nutrition factors on the
respiratory health of Pacific youth. The results suggested a positive
impact of consuming more fruit and vegetables during childhood on
respiratory health later in life. The follow-up study will continue to
explore the effects of factors from relevant domains based on the PIFS
cohort, where a new integrated model will be applied. A simulation will
be conducted to determine this model.

\subsection{Elastic Net Regularization for Vector Generalized Linear
Models: A Flexible Framework for High-Dimensional Biomedical
Data}\label{elastic-net-regularization-for-vector-generalized-linear-models-a-flexible-framework-for-high-dimensional-biomedical-data}

\emph{Wenqi Zhao}\index{Zhao, Wenqi} \emph{(University of Auckland,
Auckland, New Zealand)}

\emph{Co-authors: Dr.Thomas Yee}

\setlength{\parskip}{0.5em}

We introduce a novel implementation of elastic net regularization for
vector generalized linear models (VGLMs), capable of fitting over 100
family functions and designed to support complex, high-dimensional
modeling tasks commonly encountered in the biosciences. VGLMs extend
classical GLMs by accommodating multivariate and multi-parameter
responses, making them particularly well-suited for heterogeneous
biomedical data.

Our method integrates sparse estimation techniques—such as lasso,
ridge, and their convex combinations—into this broader modeling
framework, enhancing model interpretability and stability in
high-dimensional settings. The algorithm is implemented in the vglmnet
function within the new VGAMplus package for R. It leverages a modified
iteratively reweighted least squares (IRLS) procedure, combined with
pathwise coordinate descent, Karush-Kuhn-Tucker (KKT) condition checks,
and strong rules for variable screening to ensure computational
efficiency.

This framework supports a wide range of models beyond the exponential
family, including ordinal, categorical, and zero-inflated distributions
commonly encountered in fields such as epidemiology, genomics, and
clinical research. We illustrate the utility of our approach through
comparisons with existing tools (glmnet, ordinalNet, mpath) and apply it
to real-world datasets involving survival outcomes, count data, and
bivariate binary responses. By uniting the structural flexibility of
VGLMs with the benefits of regularization, our method provides a
powerful and scalable solution for modern statistical modeling in the
biosciences.

\subsection{Extending Spatial Capture-Recapture with the Hawkes
Process}\label{extending-spatial-capture-recapture-with-the-hawkes-process}

\emph{Alec B. M. van Helsdingen}\index{van Helsdingen, Alec B. M.}
\emph{(University of Auckland)}

\emph{Co-authors: Charlotte M. Jones-Todd}

\setlength{\parskip}{0.5em}

Spatial capture-recapture (SCR) is a well-established method used to
estimate animal population size from animal sighting or trapping data.
Standard SCR methods assume animal movements are independent and
consequently cannot incorporate site fidelity (attachment to a
particular region) nor the temporal correlation of an animal’s
location. Recent work has sought to solve these issues by explicitly
modelling animal movement.

In this talk we propose an alternative solution for camera trapping
surveys based on a multivariate self-exciting Hawkes process. Here the
rates of detection of a given animal at a given camera are a function of
not only the location and its proximity to the animal’s activity
center, but also where and when the animal was most recently detected.

Through a mixture of Gaussian distributions, our model expects more
detections closer in space to the last detection, and reduces to SCR
when an animal is yet to be detected. This formulation, we believe,
better reflects animal behaviour because shortly after detection, we
expect to see an individual close to where it was last seen. Thus, our
model allows us to account for both site fidelity and the inherent
temporal correlation in detections that have not previously been
accounted for in SCR-type models.

In this talk, I will 1) give an overview of Self-Exciting Spatial
Capture-Recapture (SESCR) models, and 2) demonstrate the additional
inference that can be drawn from such models and 3) apply the framework
using a few case studies to compare traditional SCR and SESCR.

\section{Tuesday}\label{tuesday}

\section{Wednesday}\label{wednesday}

\section{Thursday}\label{thursday}

\section{Friday}\label{friday}

\chapter{Posters}\label{posters}

\subsection{The ``Theory of Sampling'' (ToS) - have statisticians missed
the
boat?}\label{the-theory-of-sampling-tos---have-statisticians-missed-the-boat}

\emph{Damian Collins}\index{Collins, Damian} \emph{(NSW DPIRD)}

\emph{Co-authors: Anne Harris}

\setlength{\parskip}{0.5em}

The ``Theory of Sampling'' (ToS) framework originated with chemical
engineer Pierre Gy over 50 years ago. It essentially deals with the
sampling of materials, particularly soils, composts,and minerals. It
enjoys a strong advocacy amongst chemical engineers, geologists and
similar professionals. The ToS instructs practitioners to stop and think
about the sampling procedure. It encourages them to consider all the
sources of variation in the material they wish to sample rather than
just choosing to take the most convenient sample such as a “grab
sample” from the surface. However, one of its strongest advocates
argues that, in order to save on analysis costs, only one “correct”
composite sample is necessary. This single sample, deemed to be
“representative” of the entire material, is all that is required to
characterise a material (Esbenson (2017)). Composting sampling
guidelines around the world use the ToS to justify this practice of
single sampling. For example, it has currently been recommended in the
AS4454 review in Australia (https://www.as4454review.com.au/). However,
with only one sample analysed, contributions to measurement uncertainty
from either the sampling or analysis processes, cannot be assessed. It
is concerning that despite the apparent sophistication of ToS
methodology and the breadth of its advocacy network, it perpetuates the
scientifically unsound idea of a “representative sample.” .Even the
name, “Theory of Sampling”, is a misnomer, suggesting an even
broader application than just the sampling of materials. We review the
ToS and consider implications to statistical practise and adoption.

Esbenson, K.H. (2017). Sampling: Theory and Practice. Alchemist 85 pages
3–6.
https://www.lbma.org.uk/alchemist/issue-85/sampling-theory-and-practice

\subsection{Evaluating Remote Monitoring in Automated Peritoneal
Dialysis: A Difference-in-Differences
Analysis}\label{evaluating-remote-monitoring-in-automated-peritoneal-dialysis-a-difference-in-differences-analysis}

\emph{Annie Conway}\index{Conway, Annie} \emph{(ANZORRG , ANZDATA ,
Australia and New Zealand Dialysis and Transplant Registry)}

\emph{Co-authors: Stephen McDonald}

\setlength{\parskip}{0.5em}

During a cholera outbreak in 1850s London, the physician John Snow
famously identified the source of the outbreak as the water supply when
he removed the handle from one of the city's water pumps. He used a
rudimentary version of the method \emph{difference-in-differences} (DiD)
whereby both treated/untreated and before/after are compared to estimate
the effect of an intervention.

Recently, DiD has gained renewed popularity in health and
epidemiological sciences for testing the effectiveness of technological
or policy interventions. However, real-world applications can have
complexities such as staggered roll-out of the intervention and
heterogeneity between treated units. Additionally, these often involve
count outcomes that violate the linear model assumptions underlying
existing DiD methods.

In this project I consider the impact of a new technology for automated
peritoneal dialysis (APD) that allows patients to be remotely monitored
by clinicians as they complete dialysis at home. This new APD cycler was
rolled out gradually in PD centres across Australia and New Zealand
between 2017 and 2023. Using data from the ANZDATA registry, I assess
remote monitoring's impact across PD centres, counting the number of
deaths, peritonitis infections and technique failures, comparing before
and after the roll-out and between treated and untreated centres.

To address staggered roll-out and small count data, I develop a
two-stage extension of Callaway and Sant'Anna's (2021) method, using
generalised linear models to generate predicted rates, then applying
staggered DiD with bootstrapped standard errors.

\subsection{Using atmospheric transport models to predict species
incursions in northern
Australia}\label{using-atmospheric-transport-models-to-predict-species-incursions-in-northern-australia}

\emph{Zhenhua (Iris) Hao Dr}\index{Hao Dr, Zhenhua (Iris)}
\emph{(Department of Agriculture, Fisheries and Forestry)}

\setlength{\parskip}{0.5em}

Predicting the arrival of invasive species is critical for strengthening
biosecurity, particularly in regions like northern Australia that are
vulnerable to aerial incursions. This ongoing study compares atmospheric
modelling with detection patterns from long-term field surveillance to
evaluate how different models can be integrated to increase confidence
in identifying likely source regions, arrival pathways and periods of
elevated risk. We apply three complementary models: HYSPLIT (Hybrid
Single-Particle Lagrangian Integrated Trajectory), FLEXPART (FLEXible
PARTicle dispersion model), and Lagrangian Coherent Structures (LCS).
These models range in complexity - from simulating basic wind
trajectories to identifying hidden transport patterns within the
atmosphere. A rich empirical foundation is provided by two decades of
fruit fly trapping data collected through the Northern Australia
Quarantine Strategy (NAQS) in the Torres Strait. By identifying likely
incursion windows for Oriental Fruit Fly (Bactrocera dorsalis) arrivals
from Papua New Guinea, we evaluate and compare model performance and
gain insights into potential dispersal mechanisms. While model
implementation and case analysis are still underway, preliminary results
suggest that combining outputs from different models improves the
detection and interpretation of likely incursion events. Future work
will focus on integrating these model outputs with additional
environmental and biological datasets to produce dynamic risk maps.
These maps aim to support NAQS surveillance planning and resource
prioritisation. This framework can also be extended to other high-risk
species to enhance predictive capacity and improve biosecurity
strategies for Australia's north.

\subsection{Cleaning Text Data with Large Language
Models}\label{cleaning-text-data-with-large-language-models}

\emph{Jiajia Li}\index{Li, Jiajia} \emph{(Australian National
University)}

\emph{Co-authors: Emi Tanaka}

\setlength{\parskip}{0.5em}

Data cleaning – particularly the standardisation and feature
extraction from unstructured text – is a time-consuming and often
manual process in data science workflows. Traditional methods, such as
approximate string matching, lack semantic understanding and still
require inefficient, manual corrections. We leverage the transformative
potential of large language models (LLMs) to automate these tedious
tasks and significantly accelerate data preparation.

We introduce the \texttt{emend} R package, which leverages LLMs to
perform various data cleaning tasks, including matching categorical
variables, reordering ordinal variables, translating languages, and
standardising dates and addresses. The package integrates with local
LLMs via Ollama as well as cloud-based APIs via the R package ellmer,
offering multiple options for using LLMs from different providers. A
comprehensive simulation study evaluated emend’s core functionality
across 13 LLMs and a baseline string-matching method on real-world
datasets of country names, institutional affiliations, and species
names. Results demonstrate high accuracy for LLM-based methods in
general text standardisation, with GPT-4o-mini and GPT-4o excelling on
country names and affiliations, respectively. The \texttt{emend} package
streamlines data cleaning in R, reducing the need for extensive manual
corrections for data analysts. This work highlights LLMs as powerful
tools for increasing efficiency of cleaning text data.

\subsection{Performance of Factor Analytic Mixed Models and
Plackett--Luce Models for Genotype Ranking in Multi-Environment
Trials}\label{performance-of-factor-analytic-mixed-models-and-plackettluce-models-for-genotype-ranking-in-multi-environment-trials}

\emph{Jiazhe Lin}\index{Lin, Jiazhe} \emph{(ANU)}

\emph{Co-authors: Emi Tanaka, Fonti Kar}

\setlength{\parskip}{0.5em}

Multi-environment trials are central to plant breeding, providing
evaluations of genotype performance across locations and time to capture
environmental variability. The resulting data are often heterogeneous
and unbalanced, motivating the use of advanced statistical models. The
Factor Analytic Linear Mixed Model (FALMM) has become a benchmark
approach, offering a parsimonious decomposition of genotype and
genotype-by-environment effects and supporting inference on average
genotype performance. Alternatively, the Plackett-Luce Model (PLM),
originally developed in psychology, allows aggregation of
environment-specific rankings into a global ranking that reflects
overall genotype performance. We propose a two-stage framework that
integrates FALMM and PLM, and through extensive simulation, evaluate
conditions under which each method provides optimal inference.

\subsection{ggincerta: An R Package for Uncertainty Visualisation with a
Layered Grammar of
Graphics}\label{ggincerta-an-r-package-for-uncertainty-visualisation-with-a-layered-grammar-of-graphics}

\emph{Xueqi Ma}\index{Ma, Xueqi} \emph{(Australian National University)}

\emph{Co-authors: Emi Tanaka, Francis KC Hui, Weihao Li, Quan Vu}

\setlength{\parskip}{0.5em}

Uncertainty is inherent in all estimates derived from data. However,
when visualising these estimates on spatial maps, uncertainty is often
overlooked, leading to decision-making based solely on point estimates.
The \texttt{Vizumap} R package addresses this gap by integrating both
estimates and their uncertainties in areal data, offering four types of
visualisations: bivariate choropleth maps, pixel maps, glyph maps, and
exceedance probability maps.

Building upon \texttt{Vizumap}, we introduce the \texttt{ggincerta} R
package, a ground-up reimplementation that recreates all of
\texttt{Vizumap}’s visualisations, but also provides enhanced
functionality, within the layered grammar of graphics framework provided
by \texttt{ggplot2}. By leveraging \texttt{ggplot2}, \texttt{ggincerta}
seamlessly integrates uncertainty visualisation into users’ existing
workflows, employing familiar syntax and conventions. This approach not
only enhances accessibility for those already comfortable with
\texttt{ggplot2}, but also exemplifies how existing systems can be
extended to new purposes, rather than developing separate, ad-hoc
solutions that require learning additional languages or interfaces.

\subsection{Small Area Estimation of Fertility Rate in the Philippines
using DHS and Census
Data}\label{small-area-estimation-of-fertility-rate-in-the-philippines-using-dhs-and-census-data}

\emph{George Manapat}\index{Manapat, George} \emph{(University of the
Philippines Diliman)}

\setlength{\parskip}{0.5em}

Accurate and timely estimation of fertility indicators at subnational
levels is essential for effective health planning, especially in
countries with geographically diverse populations and uneven data
availability. In the Philippines, while national and regional fertility
rates are routinely measured through national surveys such as the
Demographic and Health Survey (DHS), local-level estimates, e.g., at the
provincial level, remain sparse and unreliable due to limited sample
sizes. Using the data from the 2022 Philippine DHS and auxiliary
covariates from the 2020 Census of Population and Housing, small area
estimates of fertility rate at the provincial level are obtained using a
Bayesian hierarchical model accounting for spatial correlation.
Preliminary findings show substantial differences in fertility rates
between provinces. Some areas, especially in more remote or underserved
regions, have much higher fertility than the national average. The
results underscore the necessity for area-specific health programs and
policies to address the unique demographic challenges of each locality.

\subsection{Bayesian inference for sparse Gaussian copula graphical
model}\label{bayesian-inference-for-sparse-gaussian-copula-graphical-model}

\emph{Tomotaka Momozaki}\index{Momozaki, Tomotaka} \emph{(Tokyo
University of Science)}

\setlength{\parskip}{0.5em}

We consider Bayesian inference for sparse dependence structures among
variables in multivariate mixed data, which consists of both continuous
and discrete variables. Traditional graphical modeling approaches often
struggle with mixed-type data, as they typically focus on either purely
continuous or purely discrete settings. This limitation becomes
particularly problematic in high-dimensional scenarios where identifying
relevant variable relationships is crucial for interpretability and
prediction accuracy. This study proposes a novel Bayesian approach that
combines Gaussian copula graphical models with advanced shrinkage
priors, including horseshoe priors and double exponential priors. The
Gaussian copula framework enables unified modeling of mixed-type
variables by transforming them to a common latent Gaussian scale, while
shrinkage priors facilitate automatic variable selection by adaptively
shrinking irrelevant dependencies toward zero. This combination
addresses both the mixed-data challenge and the high-dimensionality
issue simultaneously. We develop an efficient posterior sampling
algorithm that maintains the positive definiteness of correlation
matrices by employing a block Gibbs sampler with hit-and-run algorithm.
Our approach utilizes rank information to handle mixed-type data within
the copula framework, enabling unified treatment of continuous and
discrete variables. The proposed method is particularly valuable for
analyzing complex datasets in genomics, economics, and social sciences,
where variables of different types naturally co-occur. We present
comprehensive evaluation of our method's effectiveness through extensive
simulation studies comparing with existing approaches and real data
analysis demonstrating practical applicability.

\subsection{Paired Comparison with Cyclic Dominance: An Extension of the
Bradley–Terry
Model}\label{paired-comparison-with-cyclic-dominance-an-extension-of-the-bradleyuxe4uxecterry-model}

\emph{Yuki Ohno}\index{Ohno, Yuki} \emph{(Tokyo University of Science)}

\emph{Co-authors: Kouji Tahata}

\setlength{\parskip}{0.5em}

The Bradley–Terry model has become a cornerstone for analyzing
paired-comparison data, providing a principled approach to estimating
``strength'' parameters without imposing a latent, continuous ordering
on teams, treatments, or evaluators. Its transitivity constraint—each
pairwise preference must be consistent with a global ranking—yields
excellent performance in many applications, including sports analytics,
psychometric scaling, and clinical preference studies. However,
empirical datasets sometimes exhibit cyclic dominance, rendering the
classical model inadequate and leading to systematic misfit.

Real-world data often display cyclic patterns, resembling a
“rock-paper-scissors” dynamic, which can violate the transitivity
constraint typically expected in such analyses. To address this issue,
we first review important applications of the Bradley-Terry model and
highlight instances where non-transitive outcomes occur. Next, we
introduce a minimal extension of the model that preserves the original
likelihood but adds a single ``cycle-strength'' parameter. This
parameter is designed to capture uniform three-way dominance among any
trio of outcomes. When the cycle-strength parameter is set to zero, the
model reverts to the classical Bradley-Terry framework. Importantly, we
demonstrate that our cycle strength corresponds precisely to the
quasi-asymmetry measure proposed by Tahata et al.~(2004), establishing a
clear quantitative relationship between model fit and deviations from
transitivity.

\subsection{Visualisation of multinomial multilevel time-series
modelling with application to current smoking
status}\label{visualisation-of-multinomial-multilevel-time-series-modelling-with-application-to-current-smoking-status}

\emph{Alice Richardson}\index{Richardson, Alice} \emph{(Australian
National University)}

\emph{Co-authors: Sumonkanti Das, Ashis Talukder, Mu Li, and Bernard
Baffour}

\setlength{\parskip}{0.5em}

In Australia, tobacco smoking has declined overall over the last 25
years, yet substantial disparities exist across age, sex and state. In
this talk we will present the results of multilevel multinomial time
series modelling to estimate trends in smoking status (current, former
and never) across small domains defined by seven age groups, two
genders, and eight states and territories from 2001–2022.

The model is developed using ex-smokers as the reference category, as
their proportion remains relatively stable over time. Statistically
significant random intercepts and random slopes for linear time trends
are identified at the state-age-sex level. The spatio-temporal extension
of the multinomial multilevel logistic model yields detailed estimates,
helping to uncover disparities in the trends of smoking decline.
Temporal random effects at the state-sex and age-sex levels also
substantially contribute to achieving numerically consistent trend
estimates.

Visualisations of these models are especially important for conveying
results in a compelling manner to health researchers and policy makers
and program implementation teams. We will discuss the process of
arriving at our visualisations using R for this particular application.

\subsection{Judgement Post-Stratification for Covariate Adjustment in
Pairwise Comparisons in Block
Designs}\label{judgement-post-stratification-for-covariate-adjustment-in-pairwise-comparisons-in-block-designs}

\emph{Sam Rogers}\index{Rogers, Sam} \emph{(University of Adelaide)}

\emph{Co-authors: Omer Ozturk and Olena Y Kravchuk}

\setlength{\parskip}{0.5em}

The randomised complete block design (RCBD) remains one of the most
widely used experimental designs in agronomy, agriculture, and
controlled-environment research. In this study, we introduce a novel
post-experimental method, judgement post-stratification, for improving
the efficiency of pairwise treatment comparisons using plot-level
covariate information. The method operates within the RCBD framework to
focus on generalised designs with multiple replications of treatments
within blocks. Using simulations based on uniformity trials, we show
that when plots can be consistently ranked according to an expected
outcome, independent of treatment assignment, this ranking can be used
post-hoc to enhance the precision of estimated treatment differences.
This approach offers a practical and flexible tool for incorporating
auxiliary information into the analysis of block designs without
compromising the integrity of the original randomisation.

\subsection{\texorpdfstring{Evaluating Diagnostic Performance via
Bayesian \(F_1\) Score Estimation without a Gold
Standard}{Evaluating Diagnostic Performance via Bayesian F\_1 Score Estimation without a Gold Standard}}\label{evaluating-diagnostic-performance-via-bayesian-f_1-score-estimation-without-a-gold-standard}

\emph{Jun Tamura}\index{Tamura, Jun} \emph{(Yokohama City University)}

\emph{Co-authors: Yusuke Saigusa, Kouji Yamamoto}

\setlength{\parskip}{0.5em}

In clinical medicine, diseases are often diagnosed with high accuracy
based on characteristic signs and symptoms. However, the gold standard
tests used to confirm diagnoses are frequently invasive, expensive, or
impractical for population-level screening. This highlights the
importance of developing and validating simple, non-invasive screening
tests. Diagnostic performance is typically evaluated using metrics such
as accuracy, sensitivity (Se), specificity (Sp), positive predictive
value (PPV), and negative predictive value (NPV). However, when disease
prevalence is low or class imbalance exists, accuracy and Sp may not
adequately reflect test performance. \(F_1\) score, which considers both
Se and PPV, has been widely adopted in such settings as a balanced
measure of diagnostic ability. A limitation of the \(F_1\) score is its
reliance on the true disease status defined by a gold standard. In some
fields, however, the disease mechanism may be unclear or the definition
of clinical signs ambiguous, making it difficult to define a reliable
gold standard. This presents challenges for evaluating diagnostic
accuracy using conventional methods. To address this issue, we propose a
Bayesian framework that enables estimation of the \(F_1\) score without
directly observing gold standard outcomes. The approach uses latent
class analysis to estimate unobserved disease status and offers a
promising method for assessing test performance in settings where the
gold standard is unavailable or uncertain

\subsection{Visualization for departures from symmetry with the
power-divergence-type measure in square contingency
tables}\label{visualization-for-departures-from-symmetry-with-the-power-divergence-type-measure-in-square-contingency-tables}

\emph{Wataru Urasaki}\index{Urasaki, Wataru} \emph{(Department of
Information Sciences, Tokyo University of Science)}

\emph{Co-authors: Tomoyuki Nakagawa, Jun Tsuchida and Kouji Tahata}

\setlength{\parskip}{0.5em}

When the row and column variables consist of the same category in a
two-way contingency table, it is called a square contingency table. Such
tables often have an association structure concentrated along the main
diagonal, making the analysis of symmetric relationships and transitions
important. Various models and measures have been proposed to analyze
these structures to understand the changes between two variables'
behavior at two-time points or cohorts. This is necessary for a detailed
investigation of individual categories and their interrelationships,
such as shifts in brand preferences. We propose a novel correspondence
analysis (CA) framework to evaluate departures from symmetry in square
contingency tables with nominal scales, using a power-divergence-type
measure. This approach ensures that well-known divergences can also be
visualized and, regardless of the divergence used, the CA plot consists
of two principal axes with equal contribution rates. Our visualization
method enables the magnitude of each category’s departure from
symmetry to be assessed by its deviation from the origin, while
asymmetric relationships between category pairs can be interpreted
through corresponding triangle areas. Importantly, the scaling of the
departures from symmetry provided by the measure is independent of
sample size, allowing for meaningful comparisons and unification of
results across different tables. This standardization supports broader
applicability in empirical studies. We also present some considerations
from the results of the analysis with actual data. Our framework thus
offers an effective tool for studying structural shifts in categorical
data.

\chapter{Code of Conduct}\label{code-of-conduct}

The International Biometrics Society Australasian Region (IBS-AR) is
dedicated to providing a welcoming and harassment-free experience for
everyone. We do not tolerate unprofessional or inappropriate conduct,
which includes discrimination and harassment, sexual or otherwise.
Conference participants should always work towards maintaining an
environment free of harassment and discrimination. At the discretion of
the IBS-AR conference organisers and/or IBS-AR council, any participants
violating these rules may be asked to leave the conference, without any
fee refund.

Participants agree to:

\begin{itemize}
\tightlist
\item
  Be considerate in speech and actions, and actively seek to acknowledge
  and respect the boundaries of fellow attendees at all times.
\item
  Refrain from demeaning, discriminatory, or harassing behaviour and
  speech.
\item
  Alert a member of the IBS-AR council or conference organisers if they
  notice a dangerous situation, someone in distress, or violation of
  this code of conduct.
\end{itemize}

Harassment is demeaning or threatening behaviour directed at an
individual or a group of people. Harassment includes, but is not limited
to,

\begin{itemize}
\tightlist
\item
  Offensive comments related to gender, gender identity and expression,
  sexual orientation, disability, mental illness, person's lifestyle
  choices and practices, physical appearance, age, race, religion (or
  lack thereof) or beliefs;
\item
  Unwelcome sexual attention and behaviour including inappropriate
  sexual propositions; unnecessary touching; verbal or written comments
  or physical actions of a sexual nature; sexually degrading words used
  to describe an individual; a display of sexually suggestive objects or
  pictures; or sexually explicit jokes; and
\item
  Threatening behaviour including bullying, threat of violence,
  deliberate intimidation, stalking or incitement of violence. Reporting
\end{itemize}

If you are being harassed by a member/guest/participant of/at the IBS-AR
Conference, or notice that someone else is being harassed, or have any
other concerns, please contact the designated conference allies in the
conference team. The current conference allies can be found on page 4.
If the person who is harassing you is one of the conference allies, they
will recuse themselves from handling your incident and will not be given
any information without your explicit consent. We will respond as
promptly as we can.



\printindex


\cleardoublepage
\newpage

\pagestyle{empty}

\ThisULCornerWallPaper{1.1}{assets/back-cover.png}


\centering



International Biometrics Society Australasian Region

\end{document}
